% concatenated from minimalramseyarxiv.tex
\documentclass[12pt]{article}
\usepackage{amssymb,comment}
\usepackage{comment}
\setlength{\parindent}{0pt} \oddsidemargin -0.1in \evensidemargin
-0.1in \topmargin -0.0in \textheight 8.5in \textwidth 6.5in
\renewcommand{\baselinestretch}{1.2}
\usepackage{amsmath,amsthm}
%\usepackage[T1,T5]{fontenc}
%\usepackage{textcomp}
\usepackage{graphicx}
\usepackage{color}
%\usepackage{cite}
%\bibliographystyle{unsrt}

\newcounter{foo}
\usepackage{titling}

\setlength{\droptitle}{-10em}

\usepackage{url}

\def\Fq{{\mathbb{F}_q}}
\def\FF{{\mathbb{F}}}
\newcommand{\PG}{\mathrm{PG}}

\newcommand{\cA}{{\cal A}}
\newcommand{\cC}{{\cal C}}
\newcommand{\cG}{{\cal G}}
\newcommand{\cH}{{\cal H}}
\newcommand{\cI}{{\cal I}}
\newcommand{\cL}{{\cal L}}
\newcommand{\cM}{{\cal M}}
\newcommand{\cP}{{\cal P}}
\newcommand{\cS}{{\cal S}}
\newcommand{\cT}{{\cal T}}
\newcommand{\cU}{{\cal U}}

\usepackage[dvipsnames]{xcolor}
\usepackage{hyperref}
\hypersetup{colorlinks = true, citecolor = NavyBlue, linkcolor = NavyBlue, urlcolor = NavyBlue}
\usepackage[backend=biber,giveninits=true,doi=false,isbn=false,url=false,eprint=false,maxbibnames=99,sortcites]{biblatex}
\usepackage{cleveref}

\newcommand{\red}[1]{\textcolor{red}{#1}}
\newfont{\blb}{msbm10 scaled\magstep1}
\newfont{\comp}{cmr12 scaled\magstep1}
\newfont{\compb}{cmr10 scaled\magstep2}
\newfont{\sbb}{cmssbx10 scaled\magstep3}
\newfont{\sbbb}{cmssbx10 scaled\magstep5}
\newfont{\sbs}{cmssbx10 scaled\magstep1}
\newtheorem{theorem}{Theorem}
\newtheorem{lemma}{Lemma}
\newtheorem{claim}[subsection]{Claim}
\newtheorem{cl}[section]{Claim}
\newtheorem{corollary}{Corollary}
\newtheorem{proposition}{Proposition}
\newtheorem{observation}{Observation}

%\newenvironment{proof}{{\bf Proof.}}{\hfill{ }\vrule height10pt width5pt depth1pt}
\newtheorem{conjecture}[foo]{Conjecture}

\theoremstyle{definition}
\newtheorem{definition}{Definition}
\newtheorem{remark}{Remark}

\newtheorem{othertheorem}{Theorem}
\renewcommand*{\theothertheorem}{\Alph{othertheorem}}
\newtheorem{othercorollary}[othertheorem]{Corollary}

\renewcommand*{\theothercorollary}{\Alph{othercorollary}}

\makeatletter
\newtheorem*{rep@theorem}{\rep@title}
\newcommand{\newreptheorem}[2]{%
	\newenvironment{rep#1}[1]{%
		\def\rep@title{#2 \ref{##1}}%
		\begin{rep@theorem}}%
		{\end{rep@theorem}}}
\makeatother

\newreptheorem{theorem}{Theorem}
\newreptheorem{corollary}{Corollary}

\Crefname{owntheorem}{Theorem}{Theorems}
\Crefname{owncorollary}{Corollary}{Corollaries}
\Crefname{repcorollary}{Corollary}{Corollaries}
\Crefname{theorem}{Theorem}{Theorems}

\addbibresource{min1bib.bib}
\renewbibmacro{in:}{} %no "in"+ journal
\DeclareFieldFormat{journaltitle}{{#1},} % italic journal title with comma
\DeclareFieldFormat[inbook,thesis]{title}{\mkbibemph{#1}\addperiod} % italic title with period
\DeclareFieldFormat[article]{title}{\mkbibemph{#1}} % title of journal article is printed as normal text
%\DeclareFieldFormat{pages}{#1} % no pp. for pages

%\newcommand{\viet}[1]{{\fontencoding{T5}\selectfont#1}}




%%%%%%%%%%%%%%%%%%%%%%%%%%%%%%%%%%%%%%%%%%%%%%%%%%%%%%%%%%%%%%%%%%%%%%%%
%%%%%%%%%%%%%%%%%%%%%%%%%%%%%%%%%%%%%%%%%%%%%%%%%%%%%%%%%%%%%%%%%%%%%%%%


\title{Improved bounds for the minimum degree of minimal multicolor Ramsey graphs}

\makeatletter
\renewcommand\@date{{%
		\vspace{-\baselineskip}%
		\large\centering
		\begin{tabular}{@{}c@{}}
			Yamaan Attwa\thanks{Institue for Mathematics, Freie Universit\"at Berlin, Germany. Funded by the Deutsche Forschungsgemeinschaft (DFG, German Research Foundation) under Germany´s Excellence Strategy – The Berlin Mathematics Research Center MATH+ (EXC-2046/1, project ID: 390685689).} 
		\end{tabular}
		\quad \begin{tabular}{@{}c@{}}
			Sam Mattheus\thanks{Department of Mathematics, Vrije Universiteit Brussel, Pleinlaan 2, 1050 Brussels, Belgium. E-mail: sam.mattheus@vub.be. Research supported by postdoctoral fellowship 1267923N from the Research Foundation Flanders (FWO).}   \end{tabular}%
		\quad   \begin{tabular}{@{}c@{}}
			Tibor Szabó\thanks{Institute  for Mathematics, Freie Universit\"at Berlin, Germany.}
		\end{tabular}
		\quad  \begin{tabular}{@{}c@{}}
			Jacques Verstraete\thanks{ Department of Mathematics, University of California, San Diego, 9500 Gilman Drive, La Jolla, CA 92093-0112, USA. E-mail: jacques@ucsd.edu. Research supported by the National Science Foundation FRG Award DMS-1952786 and NSF Award DMS-2347832.} 
		\end{tabular}
		
		\bigskip
		\bigskip
		\today
		
		
}}
\makeatother



%%%%%%%%%%%%%%%%%%%%%%%%%%%%%%%%%%%%%%%%%%%%%%%%%%%%%%%%%%%%%%%%%%%%%%%%
\begin{document}
	\maketitle
	\begin{abstract}
		%Let $H$ be any graph. A graph $G$ is {\em $r$-Ramsey-minimal} for $H$ if every $r$-coloring of the edges of
		%$H$ gives a monochromatic copy of $H$ but no proper subgraph of $G$ has this property.
		%Let $s_r(H)$ denote the smallest minimum degree of any $r$-Ramsey minimal graph for $H$.
		%We improve various previous upper bounds 
		%%due to Fox, Grinshpun, Liebenau, Person and Szabo, who showed $s_r(K_k) = O_k(r^2(\log r)^{8k^2})$, Bishnoi, Bamberg and Lesgourgues, who showed $s_r(K_k) = O(k^5r^{5/2})$, and Bishnoi and Lesgourgues, who showed $s_r(K_k) = O(k^3r^3(\log r)^3)$
		%by showing for some absolute constant $c$ that for $k, r \geq 3$,  $s_r(K_{k+1}) \leq 2^{200(k+1)} r^2 \log r$ and also 
		%\[s_r(K_{k+1}) \leq (rk)^{2 + \frac{c}{k}}[(\log k)^2 + (\log k)^{1 + \frac{c}{k}}(\log r)^{1 + \frac{c}{k}}].\]
		%This implies for $k = \Omega(\log r)$ and $k = O(1)$ that $s_r(K_{k+1}) = O((rk\log k)^2 \log r)$.
		%This is close to best possible since for all $k,r \geq 3$, $s_r(K_{k+1}) \geq k^2$ and 
		%for $k = O(1)$, $s_r(K_{k+1}) = \Omega(r^2 \log r/\log \log r)$.
		
		We provide two novel constructions of $r$ edge-disjoint $K_{k+1}$-free graphs on the same vertex set, each of which has the property that every small induced subgraph contains a complete graph on $k$ vertices. The main novelty of our argument is the combination of an algebraic and a probabilistic coloring scheme, which utilizes the beneficial algebraic and combinatorial properties of the Hermitian unital. These constructions improve on a number of upper bounds on the smallest possible minimum degree of minimal $r$-color Ramsey graphs for the clique $K_{k+1}$ when $r\geq c\frac{k}{\log^2 k}$ and $k$ is large enough.  
		%We also resolve a question of Tran regarding the colored semisaturation of complete graphs. 
	\end{abstract}
	\section{Introduction}
	%\subsection{}
	%In 1976, Burr, Erd\H{o}s and Lov\'asz~\cite{burr1976graphs} initiated a systematic study of various properties of Ramsey graphs. TSz: Are they really studying PROPERTIES or just EXTREMAL VALUES OF PARAMETERS?
	We say that a graph $G$ is {\em $r$-Ramsey} for a graph $H$, denoted by $G \to (H)_r$ if every $r$-colouring of the edges of $G$ contains a monochromatic copy of $H$.  A graph $G$ is called {\em $r$-Ramsey-minimal for $H$} if it is $r$-Ramsey for $H$, but no proper subgraph of it is. The set of all $r$-Ramsey-minimal graphs for $H$ is denoted by $\cM_r(H)$. The classical Ramsey number $R_r(H)$, one of the %best
	most well-studied parameters in Combinatorics, is then the smallest number of vertices of a graph in $\cM_r(H)$. 
	Following the pioneering work of Folkman~\cite{folkman1970graphs} on the smallest clique number of Ramsey graphs for the clique, Burr, Erd\H{o}s and Lov\'asz~\cite{burr1976graphs} in 1976 initiated the systematic study of the extremal behaviour of several other graph parameters. 
	%In~\cite{burr1976graphs} Burr, Erd\H{o}s, and Lov\'asz investigated maximum/mainimum values of different parameters for graphs in $\cM_r(H)$, 
	In their seminal paper they investigated the chromatic number, the maximum and the minimum degree, and the connectivity.
	%Since their work, many researchers have studied various properties of this class and the graphs it contains. TSz: Again are they really studying PROPERTIES?
	Subsequently, their work inspired many further investigations, e.g. \cite{horn2014degree, jiang2013degree, zhu1998chromatic, rodl2008ramsey,conlon2023three, nevsetvril1976ramsey, folkman1970graphs, fox2016minimum, han,bamberg2022minimum}. 
	%TSz: THIS IS CERTAINLY NOT TRUE FOR every H: For example, it follows from the work of R\"odl and Siggers~\cite{RS08} that $\cM_r(H)$ is infinite. TSz: I would delete as a somewhat random example.
	
	%We are particularly interested in the case when the graph parameter is the minimum degree and $H = K_{k+1}$, the complete graph on $k+1$ vertices.  
	In this paper we will be particularly interested in the minimum degree of minimal Ramsey graphs. For a graph $H$ and number $r$ of colors we define %use the following definition for any two positive integers $k$ and $r$
	\[s_r(H) := \min\{\delta(G) \,\, | \,\, G \in \cM_r(H)\},\]
	to be the smallest possible minimum degree that could occur among minimal $r$-Ramsey graphs for $H$.  For cliques in the classical two-color case, Burr et al.\ ~\cite{burr1976graphs} established the following precise result: 
	\begin{equation}
		s_2(K_{k})=(k-1)^2.
	\end{equation}
	Upon first glance, this result looks extremely surprising. 
	%for a couple of reasons.
	First, it determines the {\em exact} value of the smallest possible minimum degree in a minimal 2-Ramsey graph for $K_{k}$. This is in stark contrast with our knowledge about the smallest {\em number of vertices} in such graphs, which is hopelessly out of reach. Furthermore, the value of the smallest minimum degree turns out to be just a quadratic function of $k$. This is incredibly small considering that we know that even the smallest of Ramsey graphs will have exponentially many vertices. How could it then be possible to create a (necessarily enormous) $2$-Ramsey graph for $K_{k}$, that has a vertex with just $(k-1)^2$ neighbors, such that the presence of this vertex, and in fact any of its incident edges are {\em crucial} in guaranteeing the $2$-Ramsey-ness of said enormous graph?
	
	Fox, Grinshpun, Liebenau, Person and Szab\'o~\cite{fox2016minimum} investigated the behaviour $s_r(K_{k})$ for more than two colors. They found that for any fixed clique order $k\geq 3$ there exist positive constants $c_k,C_k$, such that 
	\begin{equation} \label{eq:fox-constantk}
		c_kr^2\frac{\log r}{\log\log r} \leq s_r(K_{k}) \leq C_kr^2(\ln r)^{8k^2}.
	\end{equation}
	For the triangle $K_3$, a slightly stronger lower bound was given in \cite{fox2016minimum}, which was proved to be tight up to a constant factor by Guo and Warnke~\cite{guo2020packing}. 
	\begin{equation}
		s_r(K_{3}) = \Theta (r^2 \log r).
	\end{equation}
	These results establish that for any fixed clique order $k$ the value of the smallest minimum degree $s_r(K_k)$ is quadratic in the number $r$ of colors, up to some logarithmic factor. The power of the logarithm in the gap between the upper and lower bounds however depends significantly on $k$. 
	
	On the other end of the spectrum, when the number $r$ of colors is constant, H\`an, R\"odl, and Szab\'o ~\cite{han} determined the order of magnitude of the smallest minimum degree $s_r(k)$, up to a logarithmic factor. 
	More generally, they have shown that there exists a constant $C$ such that for every $k^2 >r$
	%there exists a constant $C_r$ such that 
	%\begin{equation}\label{eq:han-constantk}
	%  s_r(K_{k}) \leq C_r k^2 \log^2 k.  
	%\end{equation}
	% Since one can see that $s_2(K_k) \leq s_r(K_k)$, coupled with the result of Burr et al.\ above, this result establishes that having more than two colors does not increase the smallest minimum degree $k^2$ by much, by at most a $C_r\log^2 k$-factor only.  
	\begin{equation} \label{eq:han-general}
		s_r(K_k) \leq C r^3k^2 \log^3 r \log^2 k.
	\end{equation}
	Considering that we know from \cite{fox2016minimum} and \cite{burr1976graphs} that
	$s_r(K_k) \geq s_2(K_k) = (k-1)^2$, the bound (\ref{eq:han-general}) establishes that $s_r(K_k)$ is quadratic up to a $\log^2$-factor for any fixed number $r$ of colors.
	
	When both $k$ and $r$ are increasing, say $r=r(k)$ is a decent increasing function of $k$, the best known upper bounds vary depending on how fast $r$ grows. In the range $r(k) < k^2$, the upper bound of (\ref{eq:han-general}) is the best we know.
	
	In the complementary range of $r(k)\geq k^2$ another construction of Fox et al.\ \cite{fox2016minimum} gives\footnote{In \cite{fox2016minimum} only the weaker upper bound $s_r(K_k) \leq q^3=O(r^3k^6)$ is stated. However, Bishnoi and Lesgourgues~\cite{bishnoi2022new} recently observed that the choice $q\sim rk^2$ used in \cite{fox2016minimum} for the parameter $q$ is suboptimal and the calculation there also works with $q \sim rk \log k$, resulting in (\ref{eq:fox-polynomial}).} 
	\begin{equation}\label{eq:fox-polynomial}
		s_r(K_k) = O(r^3k^3 \log^3 k).
	\end{equation}
	Bamberg, Bishnoi, and Lesgourgues~\cite{bamberg2022minimum} developed a generalization of this construction and used that to obtain
	\begin{equation}\label{eq:bamberg}
		s_r(K_k) = O(r^{5/2}k^5).
	\end{equation}
	This represents the best known upper bound when $r(k)\geq \Omega \left(\frac{k^4}{\log^6 k}\right)$ and $k$ tends to infinity.\\
	
	The minimum degree of minimal Ramsey graphs has also been the subject of considerable study beyond the quantitative behavior of $s_r(K_k)$: \cite{bishnoi2023minimum, clemens2015minimum} deal with the generalizations of the parameter to asymmetric settings and hypergraphs, \cite{szabo2010minimum, boyadzhiyska2025ramsey, grinshpun2017minimum} tackle the question of \textit{$q$-Ramsey simplicity}, while \cite{boyadzhiyska2022minimal} investigates the {\em number} of vertices that can attain the degree $s_r(H)$.
	
	
	\subsection{Our results} 
	
	\subsubsection{Upper bounds}
	Summarizing the above: (1) When either the order $k$ of the clique or the number $r$ of colors is constant, the smallest minimum degree $s_r(K_{k})$ is quadratic, up to poly-logarithmic factors, in terms of $r$ and $k$, respectively; (2) when $k$ and $r$ both tend to infinity, the best known upper bounds are polynomial, with the degree of $r$ (and sometimes also of $k$) being more than two.  
	Bamberg et al.\ \cite{bamberg2022minimum} in fact conjectured that an upper bound $r^2k^2$, up to logarithmic factors, should hold in all ranges of the parameters. 
	
	In this paper we give new constructions which establish this for a large range of the parameters and improve the best known upper bounds for every large enough $k$ and $r\geq c \frac{k}{\log^2 k}$.
	
	Our first main theorem improves the bounds (\ref{eq:bamberg}), (\ref{eq:fox-polynomial}) and (\ref{eq:han-general}) whenever $k$ tends to infinity and $r$ is large enough.
	\begin{theorem}\label{thm: MinDegGen}
		For all sufficiently large $k,r$ satisfying $k \leq r \log^2 r$, we have 
		\[s_r(K_{k}) \leq 2^{400} k^2r^{2+\frac{30}{k}} \log^{20}r\log^{20}k.\]
	\end{theorem}
	Note that this upper bound is of the form $(rk)^{2+o(1)}$ and the error term becomes logarithmic when $r(k) = e^{O(k\log k)}$. 
	
	For constant $k$, our second main result reduces the power of the log factor in the upper bound of (\ref{eq:fox-constantk}) from $8k^2$ to $2$.
	\begin{theorem}\label{thm: MinDegConst}
		For all $k \geq 3$ there exists a constant $C_k$ such that for all $r \geq 2$
		\[s_r(K_{k}) \leq C_k (r\log r)^2. \]
	\end{theorem}
	Combined with the lower bound of (\ref{eq:fox-constantk}), Theorem~\ref{thm: MinDegConst} determines the value of $s_r(K_k)$ up to a factor $O(\log r \log\log r)$, for every fixed $k\geq 4$. 
	
	\subsubsection{Colored semisaturation numbers}
	
	Tran~\cite{tran2022two} observed that a certain graph parameter, introduced by Dam\'asdi et al.\ ~\cite{damasdi2021saturation}, related to colored saturation, is also relevant for $s_r(K_k)$.
Given integers $r,k \geq 2$, let ${\cal RC}_r(K_k)$ be the set of edge $r$-colorings of complete graphs, for which any extension to an edge $r$-coloring of a complete graph of one larger order creates a new monochromatic copy of $K_k$. The $r$-color semisaturation number $\mathrm{ssat}_r(K_{k})$ is defined to be the {\em smallest} order $n$ such that there exists an edge $r$-coloring of $E(K_n)$ in ${\cal RC}_r(K_k)$. 
%The motivation behind this notion is that the classical $r$-color Ramsey number $R_r(K_k)$ is exactly $1$ plus the {\em largest} integer $n$ such that there exists a monochromatic $K_k$-free $r$-coloring of $E(K_n)$ in ${\cal RC}_r(K_k)$. 
%	Dam\'asdi et al.\ ~\cite{damasdi2021saturation} defined $\mathrm{sat}_r(K_{k})$ to be the {\em smallest} integer $n$ such that there exists a monochromatic $K_k$-free $r$-coloring of $E(K_n)$ in ${\cal RC}_r(K_k)$. Thus $\mathrm{sat}_r(K_{k}) < R_r(K_{k})$. 
%	This definition was inspired by the definition of Erd\H os, Hajnal and Moon~\cite{erdos1964problem} of the saturation number $\mathrm{sat}(n,H)$ opposite of the Tur\'an number $\mathrm{ex}(n,H)$. 

%Given an integer $r \geq 2$, let $\cG_r$ denote the family of complete graphs whose edges are colored with $r$ colors. For $G, G' \in \cG_r$, we say that $G'$ extends $G$ if $G'$ can be obtained from $G$ by iteratively adding a new vertex and colored edges connecting the new vertex to the existing ones. A member $G$ of $\cG_r$ is called $(r,K_{k+1})$-\textit{semisaturated} if every $G' \in \cG_r$ that extends $G$ must contain a new monochromatic copy of $K_{k+1}$. 
	%If moreover every color class of $G$ is $K_{k+1}$-free, then $G$ is said to be $(r,K_{k+1})$-saturated. Observe that the order of the largest $(r,K_{k+1})$-saturated graph is one less than the multicolor Ramsey number of $K_{k+1}$.%
%	We call the order of the smallest $(r,K_{k+1})$-semisaturated graph the {\em $r$-color semisaturation number of $K_{k+1}$} and denote it by %$\mathrm{sat}_r(K_{k+1})$ and %
%	$\mathrm{ssat}_r(K_{k+1})$. %respectively%
	This quantity was first investigated, within a more general framework, by Dam\'asdi et al.\ \cite{damasdi2021saturation}. Tran \cite{tran2022two} observed that 
	\begin{equation} \label{eq:semisat}
		\mathrm{ssat}_r(K_k) \leq s_r(K_k)
	\end{equation}
	and asked \cite[Question 4.2]{tran2022two} whether there exists a constant $C$ (independent of $k$) such that $\mathrm{ssat}_r(K_{k}) = O_k(r^2(\log r)^C)$.
	%motivated by earlier upper bounds on $s_r(K_k)$. 
	Our \Cref{thm: MinDegConst} together with (\ref{eq:semisat}) answers this question in the affirmative. 
	\begin{corollary}
		For all $k \geq 4$, there exists a constant $C_k$ such that  $\mathrm{ssat}_r(K_k) \leq C_kr^2 \log^2 r$.
	\end{corollary}
	
	It turns out however that for $\mathrm{ssat}_r(K_k)$ one can prove an even stronger bound. 
	In the same paper, Tran asks whether $\mathrm{ssat}_r(K_{k+1}) = \omega(r^2)$ as $r \to \infty$ \cite[Question 4.1]{tran2022two}. The following theorem answers this question in the negative.
	\begin{theorem}\label{thm:semisaturation}
		For all $k,r \geq 2$, we have $\mathrm{ssat}_r(K_{k}) \leq 4(k-2)^2r^2$.
	\end{theorem}
	When $k$ is fixed and $r$ goes to infinity, \Cref{thm:semisaturation} together with the lower bound of \cite{fox2016minimum} from (\ref{eq:fox-constantk}) establishes a separation by a factor $\frac{\log r}{\log\log r}$ between the orders of magnitude of $\mathrm{ssat}_r(K_{k})$ and $s_r(K_k)$.
	%  Recall that Fox et al \cite{fox2016minimum} proved that $s_r(K_k)\geq c_k r^2\ln r/\ln \ln r$ for some $c_k >0$. This alongside \Cref{thm:semisaturation} already show an asymptotic separation between $\mathrm{ssat}_r(K_{k})$ and $s_r(K_k)$ when $k$ is fixed and $r$ grows arbitrarily.
	%The main objective of this paper is to study the more general multi-color setting of the same parameter. Instead of dealing with $s_r(K_k)$ directly however, we work with an equivalent combinatorial parameter first mentioned in ~\cite{burr1976graphs}, later formalized explicitly in ~\cite{fox2016minimum}, and for which we need a couple of definitions.   
	
	
	\subsubsection{Lower bounds}
	Except for the (nearly) extremal ranges, i.e. when either the clique order $k$ or the number $r$ of colors is (nearly) constant, we do not have a lower bound on $s_r(k)$ which is quadratic both in $k$ and $r$ up to poly-logarithmic factors. 
	Damásdi et al.\ ~\cite{damasdi2021saturation} gave a lower bound of %{\bf Please check this:} $(r-1)k^2 - 3rk +4k +2r-3 =  
    $\Omega (k^2r)$ on $\mathrm{ssat}_r(K_k)$ that transfers to a lower bound on $s_r(K_k)$ via the observation (\ref{eq:semisat}) of Tran~\cite{tran2022two}.
	%{\bf Is your bound better?}
    %\red{To use that paper we use $c = r$ and $k_1 = \dots = k_c = k$ and hence I find $(k-1)((r-1)k-2r+3)$ which is not quite the same right?}
	
	Here we prove a lower bound that is quadratic in $r$ and linear in $k$. 
	
	\begin{theorem}\label{thm: LowerKk+1}
		For all $k,r \geq 3$, we have $s_r(K_{k}) = \Omega(kr^2)$.
	\end{theorem}
	
	
	\section{On the proof}
	Instead of dealing with minimal Ramsey graphs, the proofs of all the known bounds on $s_r(K_k)$ work with an alternative function, distilled by Fox et al.\ \cite{fox2016minimum} from the original argument of Burr et al.\ \cite{burr1976graphs} for $s_2(K_{k}) =(k-1)^2$. 
	
	\begin{definition}[Color Pattern.]
		A sequence of pairwise edge-disjoint graphs $G_1,\dots,G_r$ on the same vertex set $V$ is called an \textit{$r$-color pattern} on $V$ (where the edges of $G_i$ are said to have color $i$). The color pattern is \textit{$K_{k+1}$-free} if $G_i$ is $K_{k+1}$-free for every $i=1, \ldots r$.\\  
		Given a color pattern $G_1,\dots,G_r$ on the vertex set $V$ and an $r$-coloring $c:V\rightarrow [r]$ of the vertices, a \emph{strongly monochromatic} copy of a graph $H$ according to $c$ is a copy of $H$ whose edges and vertices all have the same color.
	\end{definition}
	
	\begin{definition}
		Let $r,k \geq 2$ be positive integers, we define $P_r(k)$ to be the smallest positive integer $n$ such that there exists a $K_{k+1}$\textit{-free} color pattern $G_1,\dots,G_r$ on the vertex set $[n]$ such that every $r$\textit{-coloring} of $[n]$ induces a strongly monochromatic $K_k$.
	\end{definition}
	The connection between $s_r$ and $P_r$ is summarized in the following lemma.
	\begin{lemma}\cite[Theorem 1.5]{fox2016minimum} \label{lem: ColorPattern}
		For all integers $r,k \geq 2$ we have
		$s_r(K_{k+1})=P_r(k).$
	\end{lemma}
	
	%Tibor: It is NOT because we want to ''ease." It is because we have NO OTHER IDEA. 
	%To ease the study of $P_r(k)$, 
	
	To prove an upper bound on $P_r(k)$, one needs to construct a $K_{k+1}$-free $r$-color pattern $G_1, \ldots , G_r$ with the specific property about strongly monochromatic $K_k$. 
	As it happens, at the moment we have no other idea of guaranteeing the existence of a strongly monochromatic $K_k$ in an arbitrary $r$-coloring of the vertices, but requiring that {\em each} of the graphs $G_i$ has a $K_k$ in {\em each} subset of size at least $n/r$ and then use this for the largest color class in $[n]$. 
	To this end, we define for a graph $G$ and a positive integer $k$ the parameter $\alpha_k(G)$ to be the order of the largest $K_k$\textit{-free} induced subgraph of $G$. 
	
	\begin{observation}\cite[Lemma 4.1]{fox2016minimum} \label{obs: ColorPatter}
		If there exists a $K_{k+1}$-free color pattern $G_1, \ldots , G_r$ on $[n]$ such that $\alpha_k(G_i) <\frac{n}{r}$ for every $i=1, \ldots , r$, then 
		$s_r(K_{k+1})=P_k(r) \leq n$.
	\end{observation}
	
	We will also follow this road and construct $K_{k+1}$-free $r$-color patterns on $[n]$ with $\alpha_k$-values less than $n/r$. 
	Before constructing $r$ edge-disjoint $K_{k+1}$-free graphs with $\alpha_k < n/r$ however, one better deals with the ''simpler'' problem of constructing just one. 
	This is exactly the task of the well-studied Erd\H os-Rogers function $f_{k,k+1}(n)$ which asks for the smallest value of $\alpha_k(G)$ of $K_{k+1}$-free graphs $G$ on $n$ vertices. 
	Given a good Erd\H os-Rogers graph, one then ''only'' has to pack as many of them as possible on $n$ vertices.  
	Indeed, Fox et al.\ ~\cite[Conjecture 5.2]{fox2016minimum} even predicted that for every fixed $k\geq 3$ we will have $P_r(k) = \Theta (r \cdot (f_{k,k+1}(r))^2)$. Considering the recent improvements of Mubayi and Verstraete~\cite{mubayi2024order} on the Erd\H os-Rogers function, Theorem~\ref{thm: MinDegConst} comes within a log-factor of resolving this conjecture.
%will also follow this road and construct $K_{k+1}$-free $r$-color patterns on $[n]$ with $\alpha_k$-values less than $n/r$. 
%	A ''simpler'' problem closely related to this is the following: For which $n$ can we construct only a single graph $G_1$ with $\alpha_k(G_1) < n/r$? This is exactly the problem of the well-studied Erd\H os-Rogers function $f_{k,k+1}(n)$ which asks for the smallest value of $\alpha_k(G)$ of $K_k$-free graphs on $[n]$. 
%	Given a good Erd\H os-Rogers graph, one ''only'' has to tile a large fraction of the complete graph with it. Indeed, Fox et al.\ ~\cite[Conjecture 5.2]{fox2016minimum} even predicted that for every fixed $k\geq 3$ we will have $s_r(K_k) = \Theta (r \cdot (f_{k-1,k}(r))^2)$. Considering the recent improvements of Mubayi and Verstraete~\cite{mubayi2024order} on the Erd\H os-Rogers function, Theorem~\ref{thm: MinDegConst} comes within a log-factor of resolving this. 
	Good constructions for the Erd\H os-Rogers function will be extremely useful for us as well, but creating color patterns using them requires additional ideas. 
    
%	The generic Erd\H os-Rodgers construction starts with an appropriate linear hypergraph on $n$ vertices and creates a graph by dropping a random $K_{k+1}$-free graph on each of its hyperedges. 
%        This random choice has to balance that no $K_{k+1}$ is created from the edges coming from within different hyperedges, but there are enough edges so that any $n/r$-subset of the vertices contains a $K_k$. 
	
	The general approach to construct the desired color patterns is to start with an ''appropriate'' $u$-uniform linear hypergraph on $[n]$ with essentially as many hyperedges as possible, that is, in the order $\frac{n^2}{u^2}$. Then one assigns one of $r$ colors to each hyperedge ''appropriately'', to indicate which color pairs of vertices inside the hyperedge will receive should they be chosen to be an edge at all. Here the linearity of the hypergraph plays a crucial role: every pair of vertices belongs to (at most) one hyperedge. Finally, one constructs the graphs $G_i$ by dropping an ''appropriately'' random $k$-partite graph within each hyperedge of assigned color $i$, where the choices for different hyperedges are usually independent. 
    This random choice must balance that no $K_{k+1}$ is created from the edges coming from within different hyperedges, yet there are enough edges so that any $n/r$-subset of the vertices contains a $K_k$. The crux of the matter is how to define the various occurrences of "appropriate" above so that they complement each other well.
	
	In the construction of Fox et al.\ ~\cite{fox2016minimum} for (\ref{eq:fox-constantk}) (crucially making use of the Erd\H os-Rogers construction of Dudek, Retter, and R\"odl~\cite{dudek2014generalized})
	and that of H\`an et al.\ ~\cite{han} for (\ref{eq:han-general}) the linear hypergraph is essentially given by the lines of an (affine or projective) plane of order $q$ and the ''appropriate'' color assignment chosen randomly. In the constructions of Fox et al.\ ~\cite{fox2016minimum} for (\ref{eq:fox-polynomial}) and of Bamberg et al.\ ~\cite{bamberg2022minimum} for (\ref{eq:bamberg}) the linear hypergraph is given by some (pseudo)lines in a higher dimensional space and the color-assignment is defined algebraically. The point of these assignments is to ensure that the hypergraph of each color class is {\em triangle-free}, hence the $K_{k+1}$-freeness of each $G_i$ will be automatic once the graphs inside the hyperedges are $K_{k+1}$-free.
	
	In our construction we also start with the projective plane, working in the dual setup, so the lines will correspond to vertices and the vertices correspond to the hyperedges. We choose the order to be $q^2$, so we are able to make use of Hermitian unitals and its beneficial algebraic and combinatorial properties. One of the main ideas of our construction is to combine the probabilistic color assignment to the hyperedges with an appropriate algebraic one. Unlike in \cite{fox2016minimum} and \cite{bamberg2022minimum}, our algebraic color assignment will not guarantee immediate $K_{k+1}$-freeness, but will however ensure that the analysis of $K_{k+1}$-freeness will only have to consider very limited types of forbidden events. The random part of the color assignment then helps to limit the number of bad events within those types. 
%	For the appropriate algebraic color assignment we take the order of the projective plane to be a square $q^2$ and use the properties of the Hermitian unital.
	% Not here \cite{mubayi2024order, mattheus2024asymptotics}.
	
	%we define for a graph $G$ and a positive integer $k$ the parameter $\alpha_k(G)$ to be the order of the largest $K_k$\textit{-free} induced subgraph of $G$. Furthermore, we say that a $K_{k+1}$\textit{-free} graph $G$ on $n$ vertices is $(n,\alpha,k)$\textit{-critical} if $\alpha_k(G)< \alpha$. We note the following implication: 
	%Tibor: I would avoid the use of the critical graphs. Just another non-descriptive definition to remember.
	%\[\text{$\exists$ $r$ edge disjoint $(n,n/r,k)\textit{-critical}$ graphs packed on the same $n$ vertices} \implies P_r(k)\leq n.\]
	%Indeed any $r$\textit{-coloring} of the vertex set contains one color class of size at least $n/r$ which in turn contains a $K_k$ in the respective color graph. We are now ready to state our main results.
	% Tibor: I would not state these here, with the same long formula repeated.
	%\begin{theorem}\label{thm:ConstantKk}
	%    For every $k \geq 3$ there exists a constant $C_k$ such that for all $r \geq 2$, there exists an integer $n \leq C_k(r \log r)^2$ and $r$ edge-disjoint $(n,n/r,k)$\textit{-critical} graphs on $[n].$  
	%\end{theorem}
	%\begin{theorem}\label{thm:GeneralK}
	%    For all sufficiently large $k,r$ satisfying $k \leq r \log^2 r$, there exists a positive integer $n \leq 2^{400} k^2r^{2+\frac{30}{k}} \log^{20}r\log^{20}k$ and $r$ edge-disjoint $(n,n/r,k)$\textit{-critical} graphs on $[n].$
	%\end{theorem}
	%The corollaries below concerning the minimum degree in minimal Ramsey graphs for the clique follow immediately from the above theorems:
	%\begin{corollary}
	%    For all $k \geq 3$ there exists a constant $C_k$ such that for all $r \geq 2$
	%	\[s_r(K_{k+1}) \leq C_k (r\log r)^2. \]
	%\end{corollary}
	%\begin{corollary}
	%    For all sufficiently large $k,r$ satisfying $k \leq r \log^2 r$, we have 
	%   \[s_r(K_{k+1}) \leq 2^{400} k^2r^{2+\frac{30}{k}} \log^{20}r\log^{20}k.\]
	%\end{corollary}
	\iffalse
	\textbf{History and Sketch of the Proofs.}
	The objective in both theorems above is to pack many edge disjoint  $(n,r,k)$\textit{-critical} graphs on a small vertex set. The several papers written on the subject proceed along the following lines: Start with an edge-partition of a complete graph $K_n$ into cliques, which we temporarily call \textit{\textbf{atomic cliques}}. The atomic cliques are then uniformly edge-colored according to some \textit{\textbf{coloring scheme}} with $r$ colors. The desired $(n,r,k)$\textit{-critical} graphs are then subgraphs of the color classes made by a \emph{sparsification procedure} which deletes all copies of $K_{k+1}$ from the color classes while maintaining many well-spread copies of $K_k$ in each.\\
	
	In one construction of Fox et al.\ \cite{fox2016minimum}, the vertices of $K_n$ represent points of a finite affine plane while the atomic cliques correspond to affine lines. The coloring scheme is a random one in which every atomic clique receives a color uniformly at random and the  sparsification procedure is a random $k$\textit{-partation} of each atomic clique on top of a random edge sampling. Fox et al.\ use this construction to prove that $P_r(k)\leq C_kr^2(\ln r)^{8k^2}$ where $C_k$ is a constant depending only on $k$. Prior to \Cref{thm:ConstantKk}, this construction provided the best bound for a fixed $k$ and an $r$ growing arbitrarily large.\\
	
	In another construction of Fox et al.\, the random coloring scheme is switched with an algebraic alternative. The main advantage of this switch is the added \textit{atomic triangle-free structure} in each color class: Every monochromatic triangle is contained fully in an atomic clique of the corresponding color and therefore all monochromatic $K_{k+1}$'s are highly structured and can be easily killed by a random $k$\textit{-partation} of each atomic clique. An analysis of this construction due to Bishnoi and Lesgoures \cite{bishnoi2022new} delivers the bound $P_r(k)= \mathcal{O}(k^3r^3\log^3 r)$. Bamberg, Bishnoi, and Lesgoures provide another construction with an atomic triangle-free coloring scheme based on a group-theoretic packing of certain subgeometries of the Hermitian generalized quadrangles. Their construction gives $P_r(k) = \mathcal{O}(k^5 r^{5/2})$. \Cref{thm:GeneralK} strengthens both of these bounds on the whole region where they were previously best.\\
	
	
	Hàn, Rödl, and Szábo \cite{han} provided a construction along the general direction. While they use the usual random coloring scheme, the sparsification procedure is somewhat different: The atomic cliques are randomly partitioned into $k+1$ parts; $k$ of which are small and constitute the active portion of the  atomic clique, and the large remainder part is called the \emph{black hole}. All edges inside each part are dropped while each black hole of an atomic clique is turned into an independent set. This construction provides $P_r(k)\leq \mathcal{O}((k\log k)^2(r\log r)^3 )$ for $k \geq \sqrt{r}$ large enough. This bound remains the best known in the region $k \geq r \log^2 r$. \\
	
	Our certificates for \Cref{thm:ConstantKk} and \Cref{thm:GeneralK} follow the same recipe. The vertices of our $K_n$ represent a subset of the lines of the projective plane while the atomic cliques correspond to the points. Our coloring scheme is twofold: Part algebraic and part random. For the algebraic part we use a pencil of Hermitian unitals to partition the points of the projective plane, effectively coloring each atomic clique with the unital to which its representative point belongs. The added benefit of this step is similar to that of the atomic triangle-freeness: Every monochromatic $K_{k+1}$ is highly structured. This coloring is further refined by a round of random coloring. At this stage, we blackbox a theorem of Mubayi and the fourth author \cite{mubayi2024order} to prove \Cref{thm:ConstantKk} and apply a sparsification technique akin to the black-holed random Turánization used earlier by Hàn, Rödl, and Szábo to prove \Cref{thm:GeneralK}. These two sparsification methods result in different dependencies on $r$ and $k$ in the two main theorems.\\
	\fi
	
    
    Hermitian unitals, which dictate the rigid structures of our monochromatic $K_{k+1}$'s, have been instrumental in recent results in Ramsey theory due to the second and fourth authors \cite{mattheus2024asymptotics}, Erd\H{o}s-Rogers functions due to Mubayi and the fourth author \cite{mubayi2024order}, recent results on the generalization of this function studied by Balogh, Chen and Luo \cite{balogh2025maximum} and Mubayi and the fourth author \cite{mubayi2024erd}, and improvements for the Erd\H{o}s-Rogers function $f_{k,k+2}(n)$ by Janzer and Sudakov \cite{janzer2025improved}. \\
	
	%The main novelty compared to the aforementioned applications is that we have to use multiple disjoint Hermitian unitals to reach the desired number of colors. 
    The organization of the paper is the following. In \Cref{sec:pencilunitals} we discuss the algebraic content of our color assignment, which is based on the use of multiple disjoint Hermitian unitals. 
	The probabilistic refining of the algebraic coloring is the subject of \Cref{Coloring II: Probability}. In \Cref{sec:graphs}, we complete the proofs of Theorem~\ref{thm: MinDegGen} and \ref{thm: MinDegConst} by arguing that our graphs are indeed likely to be $K_{k+1}$-free while maintaining a small $\alpha_k$. The proof of our lower bound in Theorem~\ref{thm: LowerKk+1} is given in \Cref{sec:lowerbounds}, while Section~\ref{section: Tran} contains the short proof of Theorem~\ref{thm:semisaturation}. Section~\ref{sec:remarks} collects a number of probelms remaining open.
    
    %, where we provide a lower bound on the parameter $s_r(K_{k+1})$; 
    %however, we currently fall short of a lower bound quadratic in both $k$ and $r$.
	
	
	%The previous proof sketch
	\iffalse
	\textbf{Sketch of the proofs of  \Cref{thm:ConstantKk,thm:GeneralK}} Our goal for both theorems is to pack many edge-disjoint $(n,n/r,k)$\textit{-critical} graphs on a small set of vertices.
	These graphs will be based on projective planes, in a dual way, in the sense that the roles of points and lines will be swapped. To be precise: we will define a graph on the lines of a projective plane and assign colors to the points. This naturally leads to an edge-coloring of the graph: the edge between two lines receives the color of the intersection point. This is the main idea behind many existing results in the literature, such as those mentioned in the previous sections. Our contribution is that we employ a finer coloring of the points of the projective plane, based on geometric ideas. In particular, we will use Hermitian unitals, which have been instrumental in recent results in Ramsey theory due to the second and fourth author \cite{mattheus2024asymptotics}, Erd\H{o}s-Rogers functions  due to Mubayi and the fourth author \cite{mubayi2024order}, recent results on the generalization of this function studied by Balogh, Chen and Luo \cite{balogh2025maximum} and Mubayi and the fourth author, and improvements for the Erd\H{o}s-Rogers function $f_{k,k+2}(n)$ by Janzer and Sudakov \cite{janzer2025improved}. \\
	
	The main novelty compared to the previous results is that we have to use multiple disjoint Hermitian unitals to reach the desired number of colors. This can be achieved through a well-known geometric construction and is the content of \Cref{sec:pencilunitals}. In the second step we not only have to make sure that small subsets induce a copy of $K_k$, but crucially we have to make sure that every color induces a $K_{k+1}$-free graph. This is the main technical step of \Cref{sec:graphs}. Making each color $K_{k+1}$-free can be achieved in two slightly different ways using randomness, leading to the different dependencies on $k$ and $q$. In \cref{sec:lowerbounds}, we provide two lower bounds on the parameter $s_r(K_{k+1})$; however, we currently fall short of a lower bound quadratic in both $k$ and $r$.
	\\
	\fi
	
	
	
	
	\iffalse
	\section{Introduction}
	
	In 1976, Burr, Erd\H{o}s and Lov\'asz~\cite{burr1976graphs} initiated a systematic study of various properties of Ramsey graphs. We say that a graph $G$ is {\em $r$-Ramsey} for a graph $H$, denoted by $G \to (H)_r$ if every $r$-colouring of the edges of $G$ contains a monochromatic copy of $H$. The existence of such a graph $G$, for any number of colors and any graph $H$ is a consequence of Ramsey's classical theorem \cite{ramsey1987problem}. A graph $G$ is called {\em $r$-Ramsey-minimal} for $H$ if it is $r$-Ramsey for $H$, but no proper subgraph is. The set of all $r$-Ramsey-minimal graphs for $H$ is denoted by $\cM_r(H)$. The classical Ramsey number $R_r(H)$, the best-studied parameter in Ramsey theory, is then the smallest number of vertices of a graph in $\cM_r(H)$. In their seminal paper, Burr, Erd\H{o}s, and Lov\'asz investigated different parameters of graphs in $\cM_r(H)$, such as their chromatic number, minimum, and maximum degree. Since their work, many researchers have studied various properties of this class and the graphs it contains. For example, it follows from the work of R\"odl and Siggers~\cite{RS08} that $\cM_r(H)$ is infinite. 
	
	We are particulary interested in the case when $H = K_{k+1}$, the complete graph on $k+1$ vertices. Perhaps surprisingly Burr, Erd\H{o}s and Lov\'asz~\cite{burr1976graphs} were able to determine the minimal minimum degree of graphs in $\cM_2(K_{k+1})$.
	So for that purpose, denote this quantity as
	\[s_r(K_{k+1}) := \min\{\delta(G) \,\, | \,\, G \in \cM_r(K_{k+1})\},\]
	then we can state their result as follows.
	
	\begin{theorem}[\cite{burr1976graphs}]
		For all $k \geq 2$, $s_2(K_{k+1}) = k^2$.
	\end{theorem}
	
	Their proof essentially shows, but does not explicitly state it, that the determination of $s_2(K_{k+1})$ is equivalent to another problem on packing $K_{k+1}$-free graphs. It is this reformulation that we use to prove new upper bounds on $s_r(K_{k+1})$, taking the liberty to defer the precise relation between both to \Cref{subsec:folkman,subsec:minimalramsey}. Our main technical contributions are two new edge-disjoint packings of $K_{k+1}$-free graphs on a common set of vertices, each with slightly different dependencies on the parameters.
	
	\begin{theorem}\label{thm:main}
		Let $k \geq 3$ and $q$ a prime power such that $q \geq 2^{10}k\log q$. For all non-negative integers $s \leq q^2/(2^{12}k\log q)$ there exist edge-disjoint $K_{k+1}$-free graphs $G_1,\dots,G_s$ on the same vertex set $[N]$, $q^4/2 \leq N \leq q^4$, with the additional property that every subset of size at least $2^{100k}q^2\log q$ induces a copy of $K_k$ in each $G_i$.
	\end{theorem}
	
	\begin{theorem}\label{thm:main2}
		Let $k\geq 3$ and $q$ be a prime power such that $q^{1-5/(2k-2)} \geq 48000k\log^2q$. For all non-negative integers $s \leq q^{2-5/(2k-2)}/(96000k\log^2 q)$ there exist edge-disjoint $K_{k+1}$-free graphs $G_1,\dots,G_s$ on the same vertex set $[N]$, $q^4/2 \leq N \leq q^4$, with the additional property that every subset of size at least $12000(k\log k)(q^{2+5/(2k-2)}\log q)$ induces a copy of $K_k$ in each $G_i$.
	\end{theorem}
	
	\red{Is the requirement $k \leq q^{1-5/(2k-2)}$ much of a restriction? Since in later applications we have $rk \sim q^{2-5/(2k-2)}$, it forces $r \geq q \gg k$.}
	This remark is not clear to me. Reducing the number of vertices and increasing the number of colors are both good.
	
	\begin{remark}\label{remark:morecolors}
		From the proofs of these results, a more general construction with more colors can be obtained at the cost of reducing the number of vertices and vice versa. {\color{blue} I do not understand this remark. If we can have more colors with fewer vertices, we are winning.} The simpler statement as above suffices for the corollaries in the next sections.
	\end{remark}
	
	The implications of these two results to the function $s_r(K_{k+1})$ are the following improved upper bounds.   
	
	\begin{corollary}\label{cor:cor1}
		For all $k \geq 3$ there exists a constant $C_k$ such that for all $r \geq 2$
		\[s_r(K_{k+1}) \leq C_k (r\log r)^2. \]
	\end{corollary}
	
	\begin{corollary}\label{cor:cor2}
		There exist constants $c_1,c_2,c_3 > 0$ such that for all integers $k,r$ satisfying $k \geq 4$ and $r \geq c_1k$ we have 
		\[s_r(K_{k+1}) \leq c_2(r\log^4 r)^{2+\frac{c_3}{k}}(k\log k)^2. \]
	\end{corollary}
	
	\textbf{Structure of the paper.} In \Cref{subsec:folkman,subsec:minimalramsey,subsec:saturated} give an overview of previous results concerning $s_r(K_{k+1})$ and the related packing problem. The reader familiar with these topics can skip straight to the consequences of the above results stated in \Cref{cor:CP1,cor:CP2}. We also introduce the notion of (semi)saturated Ramsey numbers, which is a closely related problem. For the latter, we provide a construction based on projective planes, answering a question of Tran~\cite{tran2022two}. This construction will serve as a stepping stone to the more complicated constructions in the proofs of \Cref{thm:main,thm:main2}. We give a sketch of the proof at the end of \Cref{subsec:saturated} and the complete proofs in \Cref{sec:pencilunitals,sec:singlegraph}.
	
	\subsection{Color patterns and Erd\H{o}s-Rogers functions} \label{subsec:folkman}
	
	We will first make the detour to edge-disjoint packings of clique-free graphs before reviewing the known bounds on $s_r(K_{k+1})$ and our improvements on them. \\
	
	A sequence of pairwise edge-disjoint graphs $G_1,\dots,G_r$ on the same vertex set $V$ is a \textit{color pattern} on $V$ (where the edges of $G_i$ are said to have color $i$). The color pattern is \textit{$K_{k+1}$-free} if every $G_i$ is $K_{k+1}$-free.
	
	For a real $\alpha > 0$ and positive integers $k,r \geq 2$, the \textit{color pattern density Folkman number} $CP_r(\alpha,k)$, introduced by Hàn, R\"odl and Szab\'o \cite{han}, is defined as the smallest integer $N$ for which there exists a $K_{k+1}$-free color pattern $G_1,\dots,G_r$ on an $N$-element vertex set such that for every $i \in [r]$ the graph $G_i$ induces a copy of $K_k$ in any subset $U \subseteq V$ of size $\alpha|V|$. They showed the following results, which implies the existence of these numbers.
	
	\begin{theorem}[\cite{han}]\label{thm:HRS2}
		For all $r \geq 2$ there exists a constant $C_r$ such that for all $k \geq 2$
		\[CP_r(1/r,k) \leq C_r(k\log k)^2.\]
	\end{theorem}
	
	For large $k$, they can obtain a reasonable expression for $C_r$.
	
	\begin{theorem}[\cite{han}]\label{thm:HRS1}
		There exists an absolute constant $k_0$ such that for all $k > \max\{k_0,r\}$ 
		\[CP_r(1/r,k) < 80^3(r\log r)^3(k\log k)^2.\]
	\end{theorem}
	
	The function $CP_r(\alpha,k)$, and its name, are closely related to the \textit{vertex Folkman number} $F(r,k,m)$. For integers $r,k,m \geq 2$, it is defined as the minimum number $N$ such that there is a graph $G$ on $N$ vertices with the property that $G$ does not contain a $K_m$, yet any $r$-coloring of the vertices of $G$ yields a monochromatic copy of $K_k$. Note the difference with the usual Folkman number, which concerns colorings of the edges, or equivalently, is the smallest number of vertices of a graph in $\cM_r(K_k)$ with clique number at most $m-1$. Again, it may not be obvious that $F(r,k,m)$ exists (i.e.\ is finite) for any choice of $k,m,r$ with $k < m$, but this was established by Folkman \cite{Folkman70} in 1970. For the most restrictive case $m = k+1$, this follows from \Cref{thm:HRS1} as one can observe that 
	\[F(r,k,k+1) \leq CP_1(1/r,k) \leq CP_r(1/r,k). \]
	
	%Indeed, in previous work of Dudek and R\"odl \cite{DR10} it was shown that $F(r,k,k+1) = O(k^2\log^4k)$, although implicitly they showed that this was a bound on $CP_{r}(1/r,k)$.
	
	Observe that in the definition of $CP_r(\alpha,k)$, every single graph $G_i$ needs to be $K_{k+1}$-free and has the property that every subset of size $\alpha|V(G_i)|$ induces a copy of $K_k$. Focusing our attention on a single graph for a moment, this means that in order to obtain good upper bounds on $CP_r(\alpha,k)$, one has to start with a $K_{k+1}$-free graph such that every small subset of vertices induces a copy of $K_k$. This is precisely the problem introduced by Erd\H{o}s and Rogers \cite{ER62} in 1962: the \textit{Erd\H{o}s-Rogers function} $f_{k,k+1}(n)$ is defined as the maximum integer $m$ such that every $n$-vertex $K_{k+1}$-free graph has a $K_k$-free subgraph on $m$ vertices. This function, and its generalization $f_{s,t}(n)$, have been extensively studied in the literature \cite{Wolfovitz13,DRR14,JS23}. For $k = 2$ this function is intimately related to the off-diagonal Ramsey number $r(3,t)$ and it is known that its order of magnitude is $\sqrt{n\log n}$ as $n \to \infty$. However, for $k \geq 3$, less is known. A lower bound, based on a result due to Shearer \cite{Shearer}, was observed by Dudek and Mubayi and shows that
	\begin{align}\label{eq:ERlowerbound}
		f_{k,k+1}(n) = \Omega \left( \frac{\sqrt{n\log n}}{\log\log n} \right).
	\end{align}
	Upper bounds on $f_{k,k+1}(n)$ were recently improved by Mubayi and the second author \cite{mubayi2024order}, based on a combination of ideas of Wolfovitz \cite{Wolfovitz13} and Dudek, Retter, R\"odl \cite{DRR14} with a recent construction due to the authors \cite{mattheus2024asymptotics}.
	
	\begin{theorem}[\cite{mubayi2024order}]\label{thm:ERupperbound}
		For all $k \geq 3$ we have
		\[f_{k,k+1}(n) \leq 2^{100k}\sqrt{n}\log n.\]
	\end{theorem}
	
	We mention at this point that this result (and a closely related construction) will serve as the building blocks for the proofs of \Cref{thm:main,thm:main2}. \\
	
	Returning to the function $CP_r(\alpha,k)$, we can now state the first applications of \Cref{thm:main,thm:main2}. Their proofs are technical and boil down to picking the right parametrisation for $q$ and verifying several inequalities.
	
	\begin{corollary}\label{cor:CP1}
		For all $k \geq 3$ there exists a constant $C_k$ such that for all $r \geq 2$
		\[CP_r(1/r,k) \leq C_k (r\log r)^{2}.\]
	\end{corollary} 
	
	\begin{proof}
		By Chebyshev's theorem, we can find for all $k \geq 3$ and $r \geq 2$ a prime $q$ such that
		\begin{align}\label{eq:parametrisation}
			2^{50k+3/2}\sqrt{r(200k+\log r)} \geq q \geq 2^{50k+1/2}\sqrt{r(200k+\log r)}.
		\end{align}
		For this $q$, one can observe that $q \geq 2^{10}k\log q$. We can hence apply \Cref{thm:main} for $s = q^2/(2^{12}k\log q)$ and find an integer $N$ satisfying $q^4/2 \leq N \leq q^4$, and graphs $G_1,\dots,G_s$ so that setting $\alpha = 2^{100k}q^2\log q/N$ we obtain the upper bound
		\[CP_s(\alpha,k) \leq N.\]
		
		By \eqref{eq:parametrisation} we see that $N \leq q^4 \leq C_k(r\log r)^2$ for some constant $C_k$ depending on $k$. Observe that the function $CP_r(\alpha,k)$ is monotonically increasing in $r$ and decreasing in $\alpha$. Hence it suffices to show that $1/r \geq \alpha$ and $r \leq s$ to conclude
		\[CP_r(1/r,k) \leq CP_s(\alpha,k) \leq C_k(r\log r)^2.\]
		
		The inequality $r \leq 1/\alpha$ can be rewritten as $2^{100k}r \leq N/(q^2\log q)$ and hence as $N \geq q^4/2$ it suffices to show that $2^{100k+1}r \leq q^2/\log q$. Observe that $\log q \leq 200k+\log r$ follows easily from \eqref{eq:parametrisation} and so we can rewrite
		\[\frac{q^2}{\log q} \geq \frac{2^{100k+1}r(200k+\log r)}{200k+\log r} = 2^{100k+1}r, \]
		as required.
		
		Recalling the expression for $s$, it remains to show $r \leq s = q^2/(2^{12}k\log q)$. Since $q^2/\log q \geq 2^{100k+1}r$, this is clearly satisfied.
	\end{proof}
	
	With a similar proof, \Cref{thm:main2} implies the following.
	
	\begin{corollary}\label{cor:CP2}
		There exist constants $c_1,c_2,c_3 > 0$ such that for all integers $k,r$ satisfying $k \geq 4$ and $r \geq c_1k$ we have 
		\[CP_r(1/r,k) \leq c_2 (r\log^4 r)^{2+c_3/k}(k\log k)^2.\]
	\end{corollary}
	
	\red{Here I had to make a choice whether to admit some extra $\log r$ factors in the bound or to make $r$ even larger in terms of $k$, something like $r \geq k\log^c k$ for some $c$. This is necessary for the condition $k \lesssim q^{1-5/(2k-2)}$.}
	
	\red{Moreover, $k = 3$ is not included in this result, but we have Corollary 1 for that. I'm not sure if I should explicitly mention this.}
	
	\begin{proof}
		Let $k \geq 4$ and $r \geq c_1k$ for some large constant $c_1$ to be determined later. By Chebyshev's theorem, we can find a prime $q$ such that $f(r,k) \leq q \leq 2f(r,k)$, where
		\begin{align}\label{eq:parametrisation2}
			f(r,k) =(24000rk\log k)^{\frac{2k-2}{4k-9}}\cdot 10(\log r+\log k)^{\frac{4k-4}{2k-7}}.
		\end{align}
		
		For this parametrisation, one can check, using $k \geq 4$ and $r \geq c_1k$ that $\log q \leq 10(\log k+\log r)$. Furthermore, the inequality $q^{1-5/(2k-2)} \geq 48000k\log^2q$ is satisfied for sufficiently large $c_1$. \\
		
		We can hence apply \Cref{thm:main2} for $s = q^{2-5/(2k-2)}/(96000k\log^2 q)$ and find an integer $N$ satisfying $q^4/2 \leq N \leq q^4$, and graphs $G_1,\dots,G_s$ so that setting $\alpha = 12000(k\log k)(q^{2+5/(2k-2)}\log q)/N$ we obtain the upper bound
		\[CP_s(\alpha,k) \leq N.\]
		
		By \eqref{eq:parametrisation2} we see, using $r \geq c_1k$, that $N \leq q^4 \leq c_2 (r\log^4 r)^{2+c_3/k}(k\log k)^2$ for some absolute constants $c_2$ and $c_3$. We will again use the monotonicity of $CP_r(\alpha,k)$ and show that $1/r \geq \alpha$ and $r \leq s$ in order to conclude
		\[CP_r(1/r,k) \leq CP_s(\alpha,k) \leq c_2 (r\log^4 r)^{2+c_3/k}(k\log k)^2.\]
		
		The inequality $r \leq 1/\alpha$ can be rewritten as $12000(k\log k)r \leq N/(q^{2+5/(2k-2)}\log q)$ and hence as $N \geq q^4/2$ it suffices to show that $24000(k\log k)r \leq q^{2-5/(2k-2)}/\log q$. Using the previously observed upper bound $\log q \leq 10(\log r + \log k)$, it hence suffices, using \eqref{eq:parametrisation2}, that
		\[24000(k\log k)r\cdot 10(\log k + \log r) \leq f(r,k)^{\frac{4k-9}{2k-2}}.\]
		
		This equality holds for all $k \geq 4$. Finally, in order to show $r \leq s = q^{2-5/(2k-2)}/(96000k\log^2 q)$, it suffices that 
		\[r \leq \frac{f(r,k)^{\frac{4k-9}{2k-2}}}{96000k\cdot 100(\log k + \log r)^2},\]
		which is readily verified for $k \geq 4$.
	\end{proof}
	
	\subsection{Minimum degree of minimal Ramsey graphs} \label{subsec:minimalramsey}
	
	The function $CP_r(\alpha,k)$ was introduced by Hiệp, R\"odl and Szab\'o \cite{han} in order to obtain upper bounds on another function which we now define.
	A graph with colored vertices and edges is called \textit{strongly monochromatic} if all its vertices and edges have the same color. Then the \textit{$r$-color $k$-clique packing number} $P_r(k)$ is the smallest integer $N$ such that there exists a $K_{k+1}$-free color pattern $G_1,\dots,G_r$ on an $N$-element vertex set $V$ with the property that any coloring of $V$ with the colors $1,\dots,r$ yields a strongly monochromatic copy of $K_{k}$. Clearly, by looking at the largest color class we observe
	\[P_r(k) \leq CP_r(1/r,k).\]
	In fact, most proofs of upper bounds on $P_r(k)$ implicitly proceed by upper bounding the function $CP_r(1/r,k)$ instead.
	
	The $r$-color $k$-clique packing number was introduced by Fox, Grinshpun, Liebenau, Person and Szab\'o \cite{fox2016minimum} because of its relation to $s_{r}(K_{k+1})$. They showed the following correspondence between the two problems, generalizing the case $r = 2$ implicitly shown by Burr, Erd\H{o}s and Lov\'asz.
	
	\begin{theorem}[\cite{fox2016minimum}]
		For all $k,r \geq 2$ we have $s_r(K_{k+1}) = P_r(k)$.
	\end{theorem}
	
	Using this correspondence, they obtained the following bounds for $r > 2$.
	
	\begin{theorem}[\cite{fox2016minimum}]
		There exist constants $c,C > 0$ such that for all $r \geq 2$, we have 
		\[cr^2\log r \leq s_r(K_3)\leq Cr^2\log^2 r.\]
	\end{theorem}
	
	This was improved by Guo and Warnke \cite{GW20}, who determined the correct exponent of the logarithm.
	
	\begin{theorem}[\cite{GW20}]\label{thm:GW}
		There exists a constant $C$ such that for all $r \geq 2$, we have $s_r(K_3) \leq  Cr^2\log r$.
	\end{theorem}
	
	Their proof implicitly shows an upper bound on $CP_r(1/r,2)$ and proceeds by randomly packing triangle-free graphs with small independence number onto the same set of $N$ vertices, covering all but a small proportion of the edges in the process.
	
	For bigger cliques, Fox et al.\ \ used a connection to Erd\H{o}s-Rogers functions to obtain the following bounds. 
	
	\begin{theorem}[\cite{fox2016minimum}]\label{thm:foxetal}
		For all $k \geq 2$ there exist constants $c_k, C_k > 0$ such that for all $r \geq 3$,
		\[c_kr^2\frac{\log r}{\log\log r} \leq s_r(K_{k+1}) \leq C_kr^2(\log r)^{8k^2}.\] 
	\end{theorem}
	
	Indeed, they conjecture \cite[Conjecture 5.2]{Foxetal} that for every $k \geq 3$ we have $s_r(K_{k+1}) = \Theta(r\cdot f_{k,k+1}(r)^2)$. The reader can now observe the relation between the lower bound on $s_r(K_{k+1})$ given above and the lower bound on $f_{k,k+1}(n)$ in \eqref{eq:ERlowerbound}. Furthermore, the upper bound in \Cref{thm:foxetal} is related to an upper bound on $f_{k,k+1}(n)$ due to Dudek, Retter and R\"odl \cite{DRR14}, which has the same exponential dependence on $k$. Fox et al.\ \ asked \cite[Question 5.3]{Foxetal} whether there is a universal constant $C$ such that $f_{k,k+1}(n) = O(\sqrt{n}(\log n)^C)$, and if yes, whether this construction generalizes to a packing of graphs. We already saw that the answer to the first question is yes due to the recent result of Mubayi and the second author (cf.\ \Cref{thm:ERupperbound}). Combining the inequality $s_r(K_{k+1}) = P_r(k) \leq CP_r(1/r,k)$ with \Cref{cor:CP1} we can answer the second question affirmatively too and improve upper bounds on the minimum degree of minimal Ramsey graphs in the regime where $k$ is fixed and $r \to \infty$.
	
	\begin{repcorollary}{cor:cor1}
		For all $k \geq 3$ there exists a constant $C_k$ such that for all $r \geq 2$
		\[s_r(K_{k+1}) \leq C_k (r\log r)^2. \]
	\end{repcorollary}
	
	Because of the asymptotic nature of their bounds and the fact that $k$ appears in the exponent of the $\log r$ factor, Fox et al.\ \ \cite{fox2016minimum} moreover proved an upper bound which was polynomial in both $r$ and $k$ by using triangle-free partial linear spaces.
	
	\begin{theorem}[\cite{fox2016minimum}]
		For all $k,r \geq 3$, $s_r(K_{k+1}) = O(k^6r^3)$.
	\end{theorem}
	
	This approach was further refined by Bishnoi, Bamberg and Lesgourgues \cite{bamberg2022minimum}, and Bishnoi and Lesgourgues \cite{bishnoi2022new} who used generalized quadrangles to obtain improved bounds.
	
	\begin{theorem}[\cite{bamberg2022minimum}]
		For all $k,r \geq 3$, $s_r(K_{k+1}) = O(k^5r^{5/2})$
	\end{theorem}
	
	\begin{theorem}[\cite{bishnoi2022new}]
		For all $k,r \geq 3$, $s_r(K_{k+1}) = O((kr\log k)^3)$.
	\end{theorem}
	
	The bounds on $CP_r(1/r,k)$ in \Cref{thm:HRS1} by Hiệp, R\"odl and Szab\'o \cite{han} can be translated into this setting, and have a better dependence on the size of the clique.
	
	\begin{theorem}[\cite{han}]
		There exists a constant $k_0$ such that for all $k > \max\{k_0,\sqrt{r}\}$
		\[s_r(K_{k+1}) < 80^3(k\log k)^2(r\log r)^3.\]
	\end{theorem}
	
	Observe that since $s_r(K_{k+1})$ is monotone in $r$, we see that $s_r(K_{k+1}) \geq s_2(K_{k+1}) = k^2$, so this bound is already asymptotically sharp in $k$ up to powers of $\log k$. The same authors actually conjecture that $s_3(K_{k+1}) \gg k^2$ as $k \to \infty$ \cite[Conjecture 2]{han}, but this conjecture remains open. 
	
	Our results directly imply the following bounds, which improves upon all of the previous polynomial bounds when $k$ is not too small.
	
	\Cref{cor:CP2} now directly implies \Cref{cor:cor2}, which improves upon all of the previous polynomial bounds in the regime where it applies.
	
	\begin{repcorollary}{cor:cor2}
		There exist constants $c_1,c_2,c_3 > 0$ such that for all integers $k,r$ satisfying $k \geq 4$ and $r \geq c_1k$ we have 
		\[s_r(K_{k+1}) \leq c_2 (r\log^4 r)^{2+c_3/k}(k\log k)^2.\]
	\end{repcorollary}
	
	When $k = \Omega(\log r)$, this shows $s_r(K_{k+1}) = \widetilde{O}(k^2r^2)$, suppressing logarithmic factors. It is plausible that this is the correct polynomial dependence of $s_r(K_{k+1})$ on $k$ and $r$. \red{Do we want to say this?}
	
	\subsection{Semisaturated Ramsey numbers} \label{subsec:saturated}
	
	Given an integer $r \geq 2$, let $\cG_r$ denote the family of complete graphs whose edges are colored with $r$ colors. For $G, G' \in \cG_r$, we say that $G'$ extends $G$ if $G'$ can be obtained from $G$ by iteratively adding a new vertex and colored edges connecting the new vertex to the existing ones. A member $G$ of $\cG_r$ is called $(r,K_{k+1})$-\textit{semisaturated} if every $G' \in \cG_r$ that extends $G$ must contain a new monochromatic copy of $K_{k+1}$. If moreover every color class of $G$ is $K_{k+1}$-free, then $G$ is said to be $(r,K_{k+1})$-saturated. Observe that the size of the largest $(r,K_{k+1})$-saturated graph is equal to the multicolor Ramsey number of $K_{k+1}$. 
	
	We are interested in the size of the smallest $(r,K_{k+1})$-(semi)saturated graphs, which will be denoted by $\mathrm{sat}_r(K_{k+1})$ and $\mathrm{ssat}_r(K_{k+1})$ respectively. These quantities were first investigated in a more general framework by Dam\'asdi, Keszegh, Malec, Tompkins, Wang and Zamora \cite{Damasdietal21} who obtained the following results.
	
	\begin{theorem}[\cite{Damasdietal21}]
		For all $k,r \geq 2$ we have the following.
		\begin{enumerate}
			\item $\mathrm{ssat}_2(K_{k+1}) = \mathrm{sat}_2(K_{k+1}) = k^2$, and
			\item $\Omega(k^2r) = \mathrm{ssat}_r(K_{k+1}) \leq \mathrm{sat}_r(K_{k+1}) \leq (k-1)^r$.
			\item $\mathrm{ssat}_r(K_{k+1}) \leq 48k^2r^{k^2}$.
		\end{enumerate}
	\end{theorem}
	
	It was observed by Tran~\cite{tran2022two} that $\mathrm{ssat}_r(K_{k+1}) \leq P_r(k)$ follows directly from the definitions. In fact, we have the following chain of inequalities.
	
	\begin{proposition}\label{prop:saturationinequalities}
		For all $k,r \geq 2$, we have $\mathrm{ssat}_r(K_{k+1}) \leq P_r(k) \leq \mathrm{sat}_r(K_{k+1})$.	
	\end{proposition}	
	
	\begin{proof}
		For the first inequality, one can take a $K_{k+1}$-free color pattern $G_1,\dots,G_r$ on $N = P_r(k)$ vertices and color the remaining edges of $K_N$ (i.e.\ those not covered by any $G_i$) arbitrarily to obtain an edge-colored graph $G \in \cG_r$. Now consider $G' \in \cG_r$ that extends $G$ and has one vertex $v$ more. Define a vertex-coloring $\phi$ of $G$ by setting $\phi(w)$, $w \in V(G)$, to be the color of the edge $vw$ in $G'$. By definition of $P_r(k)$, we then obtain a strongly monochromatic copy of $K_k$, which in $G'$ corresponds to a copy of $K_{k+1}$ by adding the vertex $v$. 
		
		For the second inequality, consider the smallest $(r,K_{k+1})$-saturated graph $G$ on $N = \mathrm{sat}_r(K_{k+1})$ vertices. This coloring defines a $K_{k+1}$-free color pattern $G_1,\dots,G_r$ on a vertex set $V$ of size $N$ by setting $G_i$ to be the graph induced by color $i$. Moreover, any coloring $\phi$ of $V$ will induce a strongly monochromatic $K_k$: we can translate this coloring to an edge-coloring of $K_{N+1}$ by adding one extra vertex $v$ and defining the color of the edge $vw$ to be $\phi(w)$ and all other edges receive the same color as in $G$. By definition of $\mathrm{sat}_r(K_{k+1})$, this coloring will produce a copy of $K_{k+1}$ in some color $i$, which implies that color $i$ induces a strongly monochromatic copy of $K_k$ in the color pattern.
	\end{proof}
	
	Using the first of these inequalities, among other ideas, Tran~\cite{tran2022two} obtained the following results.
	
	\begin{theorem}[\cite{tran2022two}]
		\mbox{ }
		\begin{enumerate}
			\item For all integers $k,r \geq 2$ such that $k \geq 6r$, we have $\mathrm{ssat}_r(K_{k+1}) \leq 36k^2r^2$.
			
			\item For all $k \geq 2$, there exists a constant $C_k$ such that for all $r \geq 3$
			\[\frac{1}{4}r^2 \leq \mathrm{ssat}_r(K_{k+1}) \leq C_k(\log r)^{8k^2}r^2.\]
			
			\item For all $k \geq 2$ and $r \geq 3$, $\mathrm{ssat}_r(K_{k+1}) \leq 8k^3r^3$.
		\end{enumerate}	
	\end{theorem}
	
	Furthermore, he asked \cite[Question 4.2]{tran2022two} whether there exists a constant $C$ (independent of $k$) such that $\mathrm{ssat}_r(K_{k+1}) = O_k(r^2(\log r)^C)$ and whether $\mathrm{ssat}_r(K_{k+1}) = \omega(r^2)$ as $r \to \infty$ \cite[Question 4.1]{tran2022two}. Our previous results on $P_r(k)$ already answer the former question in the affirmative with $C = 2$. However we will show a stronger result which also answers the latter question. 
	
	\begin{theorem}\label{thm:semisaturation}
		For all $k,r \geq 2$, we have $\mathrm{ssat}_r(K_{k+1}) \leq 4k^2r^2$.
	\end{theorem}
	
	\begin{proof}
		By Chebyshev's theorem, we can find a prime $q$ such that $(k-1)r < q < 2(k-1)r$. Consider the affine plane $(\cP,\cL)$ of order $q$ and denote its $q+1$ parallel classes of lines by $\cL_1,\dots,\cL_{q+1}$. Define a graph $G \in \cG_{q+1}$ on the points of the affine plane as follows: the edge between two vertices $x$ and $y$ receives color $i$ if and only if the span of the corresponding points is a line in $\cL_i$. One can observe that every color class is then the union of $q$ disjoint cliques $K_q$, and any two cliques of a different color meet in exactly one vertex. \\ 
		
		We will now show that for all $i \in [q+1]$ every set of size at least $|\cP|/r$ induces a copy of $K_k$ in color $i$. This will prove the theorem, as the number of vertices of $G$ is $q^2 < 4(k-1)^2r^2 < 4k^2r^2$ and one can easily recolor the graph $G$ to obtain a member $G' \in \cG_r$ by taking unions of color classes. Formally: for any surjective function $f:[q+1]\to[r]$, we can define a graph $G' \in \cG_r$ by defining $G_j' = \cup_i G_i$, where the index $i$ runs over all $i \in [q+1]$ such that $f(i) = j$, and $G_j'$ (resp.\ $G_i$) denotes the subgraph of $G'$ (resp.\ of $G$) induced by color $j$ (resp. $i$).   \\
		
		So fix a color $i \in [q+1]$ and consider a set of points $U$ of size at least $|\cP|/r = q^2/r > (k-1)q$. Since there are exactly $q$ lines in a parallel class, we find by pigeonhole principle that there must be a line $\ell \in \cL_i$ containing $k$ points of $U$. By definition of $G$, this corresponds to a copy of $K_k$ in color $i$.
	\end{proof}
	
	This result already shows asymptotic separation between $\mathrm{ssat}_r(K_{k+1})$ and $P_r(k)$ when $k$ is fixed and $r$ tends to infinity. This is not all too surprising as the $K_{k+1}$-freeness condition in the definition of $P_r(k)$ imposes rather strong conditions on the structure of every color class. \red{we conjecture that also for every k and r large enough this holds} \\
	
	\textbf{Sketch of the proofs of \Cref{thm:main,thm:main2}.}
	This construction of \Cref{thm:semisaturation} serves as a template for the more involved constructions of \Cref{thm:main,thm:main2}. These graphs will be based on a projective planes, in a similar but dual way, in the sense that the roles of points and lines will be swapped. To be precise: we will define a graph on the lines of a projective plane and assign colors to the points. This naturally leads to an edge-coloring of the graph: the edge between two lines receives the color of the intersection point. This is the main idea behind many existing results in the literature, such as those mentioned in the previous sections. Our contribution is that we employ a finer coloring of the points of the projective plane, based on geometric ideas. In particular, we will use Hermitian unitals, which have been instrumental in recent results in Ramsey theory due to the authors \cite{mattheus2024asymptotics}, Erd\H{o}s-Rogers functions Mubayi and the second author \cite{mubayi2024order}, recent results on the generalization of this function studied by Balogh, Chen and Luo \cite{BCL24} and Mubayi and the second author, and improvements for the Erd\H{o}s-Rogers function $f_{k,k+2}(n)$ by Janzer and Sudakov \cite{JS23}. \\
	
	The main novelty compared to the previous results is that we have to use multiple disjoint Hermitian unitals to reach the desired number of colors. This can be achieved through a well-known geometrical construction and is the content of \Cref{sec:pencilunitals}. In the second step we not only have to make sure that small subsets induce a copy of $K_k$ as in \Cref{thm:semisaturation}, but crucially we have to make sure that every color induces a $K_{k+1}$-free graph. This is the main technical step of \Cref{sec:singlegraph}. Making each color $K_{k+1}$-free can be achieved in two slightly different ways using randomness, leading to the different dependencies on $k$ and $q$ in \Cref{thm:main,thm:main2}.
	
	%Our proofs are based on a combination of ideas from \cite{han} and the construction due to the authors for the off-diagonal Ramsey number $r(4,t)$ \cite{mattheus2024asymptotics}. This latter graph was also the basis for other recent and related results such as the already mentioned improvement for the Erd\H{o}s-Rogers function $f_{k,k+1}(n)$ by Mubayi and the second author \cite{mubayi2024order}, recent results on the generalization of this function studied by Balogh, Chen and Luo \cite{BCL24} and Mubayi and the second author, and improvements for the Erd\H{o}s-Rogers function $f_{k,k+2}(n)$ by Janzer and Sudakov \cite{JS23}. \\
	%
	%In short, in order to prove an upper bound for the function $CP_r(1/r,n)$, we will in \Cref{sec:pencilunitals} first assign colors to the edges of $K_N$, based on finite geometrical considerations. Then we will construct the color pattern $G_1,\dots,G_r$ by letting each $G_i$ be a subgraph of the graph induced by color $i$, which ensures the edge-disjointness. The individual $K_{k+1}$-free graphs $G_i$ will be constructed in a uniform way, applying similar methods used in \cite{han,mattheus2024asymptotics,mubayi2024order}. We will show that every subset of size $m$, for some $m$ in terms of $N$, $k$ and $r$, induces a $K_k$ in every $G_i$ by a relatively simple analysis, which has better dependence on $k$ (but worse on $r$) compared to \cite{mubayi2024order}. This is the content of \Cref{sec:singlegraph}. Finally, choosing the parametrisation so that $m \geq N/r$, we obtain that $CP_r(1/r,k) \leq N \leq rm$ and prove Theorems~\ref{thm:main1} and \ref{thm:main2}.
	\fi
	\section{Coloring I: Finite Geometry}\label{sec:pencilunitals}
	We start by describing a classical technique to partition the points of the projective plane $\PG(2,q^2)$ based on a pencil of $q$ Hermitian unitals sharing a common tangent line. Recall that a Hermitian unital $U$ is defined by a non-singular Hermitian matrix $A$ over $\mathbb{F}_{q^2}$, i.e. $A^q = A^\top$, so that $H:x^\top A x^q = 0$. For example, the matrix 
	\[A = \begin{pmatrix} 1 & 0 & 0 \\ 0 & 0 & 1 \\ 0 & 1 & 0 \end{pmatrix}\]
	defines the Hermitian unital with equation $X^{q+1}+YZ^q+Y^qZ=0$. It is well-known, see for example \cite{BarwickEbert}, that every line intersects a Hermitian unital in either $1$ or $q+1$ points. We will say that a line is \textit{tangent} or \textit{secant} respectively. Some more combinatorial properties we need are the following. We refer to the book by Barwick and Ebert \cite{BarwickEbert} for these and many more properties of Hermitian unitals.
	
	\begin{lemma}\label{lem:unitalcombinatorics}
		Let $U$ be a Hermitian unital in $\PG(2,q^2)$, then we have the following properties:
		\begin{enumerate}
			\item $|U| = q^3+1$,
			\item there are $q^4-q^3+q^2$ secant and $q^3+1$ tangent lines to $U$,
			\item each point on the unital is incident to a unique tangent line and $q^2$ secant lines, and 
			\item \emph{~\cite{mubayi2024order}} for every $k\geq 3$ secants pairwise intersecting in $U$, there exists a point in $U$ incident to at least $k-1$ of them.
		\end{enumerate}
	\end{lemma}
	
	The idea is that given any Hermitian unital $U$ and a tangent line $\ell_\infty$ at the point $p_\infty$, we can consider the defining equations of both and define (with some abuse of notation) the pencil 
	\[\cU = \{U_\lambda:U+\lambda\cdot(\ell_\infty)^{q+1} = 0 \,\, | \,\, \lambda \in \Fq\}\] 
	We will see that each $\lambda \in \Fq$ defines a unital $U_\lambda$, where $U_0 = U$. Moreover, an elementary calculation will show that every line in $\PG(2,q^2)$ not through $p_\infty$ is tangent to exactly one unital in $\cU$ and secant to the $q-1$ others.
	
	These properties can be derived purely geometrically, but for the sake of concreteness, we will use coordinates. So consider $U:X^{q+1}+YZ^q+Y^qZ=0$, then it is easy to check that $Z = 0$ is a tangent line at the point $(0,1,0)$. We will denote them as $\ell_\infty$ and $p_\infty$.
	
	\begin{lemma}\label{lem:unitalpencil}
		Consider the set $\cU = \{U_\lambda\}_{\lambda \in \Fq}$ of unitals in $\PG(2,q^2)$ defined by 
		\[U_\lambda:X^{q+1}+YZ^q+Y^qZ+\lambda Z^{q+1} = 0.\]
		Then 
		\begin{enumerate}
			\item $\bigcup_{\lambda \in \Fq}\left(U_\lambda \setminus p_\infty\right) \cup \ell_\infty$ is a partition of the points of $\PG(2,q^2)$.
			\item Every line in $\PG(2,q^2)$ not through $p_\infty$ is tangent to exactly one unital and secant to all others.
		\end{enumerate}
	\end{lemma}
	
	\begin{proof}
		Every $U_\lambda$ is defined by the non-singular Hermitian matrix
		\[\begin{pmatrix}
			1 & 0 & 0 \\ 0 & 0 & 1 \\ 0 & 1 & \lambda
		\end{pmatrix}\]
		and is hence a Hermitian unital.
		
		Observe that $p_\infty \in U_\lambda$ for all $\lambda \in \Fq$ and this is the only common point of any two unitals in the pencil. Now given a point not on $\ell_\infty$ with homogeneous coordinates $(x,y,z)$, so that $z \neq 0$, it is clear that both $a:=x^{q+1}+yz^{q}+y^qz$ and $b:=z^{q+1} \neq 0$ are elements of $\Fq$, so that there is exactly one solution in $\Fq$ to the equation $a+\lambda b = 0$. This implies that every point not on $\ell_\infty$ is contained in exactly one unital of the pencil.
		
		Finally, any line $\ell$ not through $p_\infty$ intersects every $U_\lambda$ in either $1$ or $q+1$ as each of them is a Hermitian unital. So let $t$ and $s$ be the number of times that each case occurs. Since there are $q^2$ points on $\ell$ not on $\ell_\infty$, we see
		\begin{align}
			\begin{cases}
				t+s=q \\
				t+s(q+1) = q^2,
			\end{cases}
		\end{align}
		and hence $t = 1$, $s = q-1$.
	\end{proof}
	
	\begin{corollary}\label{lem:countcommonsecants}
		Let $\Lambda \subset \Fq$, then there are $q^4-|\Lambda|q^3+q^2$ common secants to the set of unitals $\{U_\lambda\}_{\lambda \in \Lambda}$.
	\end{corollary}
	\begin{proof}
		Denote by $\cT_\lambda$ the set of tangent lines to $U_\lambda$. By \Cref{lem:unitalpencil} we see that for distinct $\lambda,\lambda' \in \Lambda$ we have $\cT_\lambda \cap \cT_{\lambda'} = \{\ell_\infty\}$. We know that $|\cT_\lambda| = q^3+1$ by \Cref{lem:unitalcombinatorics} and hence $|\cup_{\lambda \in \Lambda}\cT_\lambda| = |\Lambda|q^3+1$ so that there are $(q^4+q^2+1)-(|\Lambda|q^3+1)$ lines which are secant to every $U_\lambda$, $\lambda \in \Lambda$.
	\end{proof}
	We hereby fix for every prime power $q$ an arbitrary subset $\Lambda:= \Lambda(q) \subset \mathbb{F}_q$ of size $\lfloor \frac{q}{2}\rfloor$.  We denote by $P=P(q):= \underset{\lambda \in \Lambda}{\cup} U_\lambda -\{p_\infty\}$ the union of the unitals indexed by $\Lambda$ except for $p_\infty$, by $L=L(q)$ the set of their common secants, and by $\mathcal{P}=\mathcal{P}(q):= \{U_\lambda\}_{\lambda \in \Lambda}$ the $\Lambda$\textit{-restricted} pencil. Lemma \ref{lem:countcommonsecants} and \Cref{lem:countcommonsecants} assert that
	
	\begin{equation*}
		q^4/2-q^3\leq |\Lambda|q^3 = |P| \leq q^4/2 \quad \text {     and       } \quad q^4/2 \leq |L| \leq q^4.
	\end{equation*}
	Later on we will construct  $K_{k+1}$\textit{-free} graphs using the disjoint unitals $\{U_\lambda\}_\lambda$ of $\mathcal{P}$. Before we do so, we need to partition each single unital, only this time the partition is done probabilistically. 
	
	
	\section{Coloring II: Probability} \label{Coloring II: Probability}
	The starting point here is the $\Lambda$\textit{-restricted} pencil $\mathcal{P}$. Inside each $U_\lambda$, we color its points uniformly at random with $c$ colors and repeat this for all $U_\lambda$, $\lambda \in \Lambda$, using a different set of $c$ colors for each unital. In this way we obtain a coloring of $P$ with $|\Lambda|c$ colors, which is well-defined since the unitals $\{U_\lambda\}_{\lambda \in \Lambda}$ are disjoint, except for the point $p_{\infty}$ which we exclude permanently. For each color $i$, let $P_i$ be the set of points with color $i$. We first establish a simple lemma for which we need the Chernoff bound.
	
	\begin{proposition}[Chernoff bound]
		Let $Z$ be a binomial random variable with mean $\mu$. Then for any real $\varepsilon \in [0,1]$,
		\begin{align*}
			&\mathrm{Pr}(Z > (1+\varepsilon)\mu) \leq \exp\left(\frac{-\varepsilon^2\mu}{4}\right) \,\,\,\text{ and } \\
			&\mathrm{Pr}(Z < (1-\varepsilon)\mu) \leq \exp\left(\frac{-\varepsilon^2\mu}{2}\right).
		\end{align*}
	\end{proposition} 
	
	\begin{lemma}\label{lem:coloringpoints}
		For any integer $c \leq q/(48\log q)$, there exists a $|\Lambda|c$\textit{-coloring} of $P$ such that for all colors $i$ the following holds.
		\begin{center}
			\begin{tabular}{lp{5in}}
				$(1)$ & The set $P_i$ has size at most $2 q^3/c$ and at least $q^3/2c$. \\
				$(2)$ & Every line in $L$ contains at least $q/2c$ and at most $2q/c$ points of $P_i$. \\
			\end{tabular}
		\end{center}
	\end{lemma}
	
	\begin{proof}
		Suppose that $P_i \subseteq U_\lambda$ for some $\lambda \in \Lambda$. The probability that a given point of $U_\lambda \setminus \{p_\infty\}$ receives color $i$ is exactly $1/c$, independently of the other points. Let $A_i$ be the event that $|P_i| \geq 2q^3/c$ or $|P_i| \leq q^3/2c$. The Chernoff bound then shows with $\varepsilon = 1$ and $\varepsilon = 1/2$ respectively that
		\[\mathrm{Pr}(A_i) = \mathrm{Pr}(|P_i| \geq 2q^3/c) + \mathrm{Pr}(|P_i| \leq q^3/2c) \leq \exp\left(-\frac{q^3}{4c}\right) + \exp\left(-\frac{q^3}{8c}\right) \leq 2\exp\left(-\frac{q^3}{4c}\right).\]
		It follows by the union bound that the probability of the event $\cup_i A_i$ is at most
		\[c|\Lambda|\left(2\exp\left(-\frac{q^3}{4c}\right)\right) \leq \frac{q^2}{48\log q}\exp\left(-\frac{12q^2}{\log q}\right) < \frac{1}{2}, \]
		and hence with probability more than $1/2$, none of the events $A_i$ occur.
		
		Similarly, given a line $\ell \in L$ and a color $i$, we find that
		\begin{align}
			\mathrm{Pr}(|\ell \cap P_i| < q/2c) &\leq \exp\left(-\frac{q}{8c}\right) \\
			\mathrm{Pr}(|\ell \cap P_i| > 2q/c) &\leq \exp\left(-\frac{(q-1)^2(q+1)}{4c(q+1)^2}  \right) \leq \exp\left(-\frac{q}{8c} \right).
		\end{align}
		%\red{where does the expression in the second line come from, is it not just $\mu = (q+1)/c$? as $P_i \subset U_\lambda$ for some $\lambda \in \Lambda$ and hence $\ell$ intersects it in $q+1$ points.}
        %\textcolor{blue}{YA: It was ugly and inaccurate, now I changed it so that it's just ugly. I write the bound in terms of $q/c$ simply because I dont want to write $(q+1)/c$ for the remainder of the paper.}
        %\red{The latter I definitely agree with.}
       

        
		We now use the union bound to conclude that the probability there exists a line in $L$ whose intersection with some color class is abnormal is at most
		\[2|L|\cdot c|\Lambda|\cdot \exp\left(-\frac{q}{8c}\right) \leq \frac{q^6}{48\log q}\exp(-6\log q) < \frac{1}{2},\]
		and therefore, with non-zero probability, there is a coloring which satisfies the required properties.
	\end{proof}
	
	For every $c \leq q/(48\log q)$ we fix a coloring as guaranteed by the preceding lemma. We then define for each color $i$ of the $|\Lambda|c$ colors the graph $\tilde{G}_i$ to be the graph whose vertex set is $L$ and whose edge set consists of pairs $\{\ell,\ell'\}$ such that $\varnothing \neq\ell \cap \ell' \subset P_i$. The graphs $\{\tilde{G}_i\}_i$ certify the following lemma:
	\begin{lemma} \label{lem: graphlemma}
		Let $q$ be a prime power, $c\leq \frac{q}{48\log q}$ be an integer, and $m :=c\lfloor\frac{q}{2}\rfloor$. There exists an integer $q^4/2\leq n \leq q^4$ and edge-disjoint graphs $\tilde{G}_1,\dots,\tilde{G}_m$ on the vertex set $[n]=L$ where the following  is true for every $i \in [m]$ and $k\geq 3$:
		\begin{enumerate}
			\item $\tilde{G}_i$ is the union of of edge-disjoint maximal cliques called \em{point-cliques}.
			\item The number of point-cliques is at least  $q^3/2c$ and at most $2q^3/c$.
			\item Every vertex $\ell \in L$ is a member in at least $q/2c$ and at most $2q/c$ point-cliques.
			\item For every $K_{k+1}$ in $\tilde{G}_i$, there exists a point-clique containing either exactly $k$ or exactly $k+1$ points of the clique. We call cliques of the former type $(k+1)$\textit{-fans} and from the latter degenerate.
			\item Every edge $e \in E(\tilde{G}_i)$ is contained in at most $2(2q/c)^k$ many $(k+1)$\textit{-fans}.
		\end{enumerate}
	\end{lemma}
	\begin{proof} The first four points are immediate from the previous discussion.  For the last point, consider an edge $e \in E(\tilde{G}_i).$ The edge $e$ corresponds to two secants $\ell_1,\ell_2\in L$ such that $\varnothing\neq\ell_1\cap\ell_2\subset P_i.$   There are two types of $(k+1)$\textit{-fans} in $\tilde{G}_i$ containing both $\ell_1,\ell_2$: Fans with $\ell_1\cap\ell_2$ as their concurrence point, of which there are at most $(|\ell_1|-1)(|\ell_2|-1)\binom{2q/c}{k-2}$; and fans with other concurrence point, of which there are at most $(|\ell_1|+|\ell_2|-2)\binom{2q/c}{k-1}$.
	\end{proof}
	\begin{remark}\label{Rem:MVConstant}
		Mubayi and the fourth author \cite{mubayi2024order} proved that for every $a\geq 128$ and sufficiently large $q$ so that $q \geq a \log q$, any graph on $n \in [q^4/2,q^4]$ vertices satisfying $(1)-(5)$ with $c=\lceil\frac{q}{a \log q}\rceil$ contains a $K_{k+1}$\textit{-free} spanning subgraph $H$ such that $\alpha_k(H)\leq 2^{40k+1}aq^2\log q$.
	\end{remark}
	
	\section{Proof of Main Theorems}\label{sec:graphs}
	We are now ready to prove \Cref{thm: MinDegGen,thm: MinDegConst}. The proof of \Cref{thm: MinDegConst} is now straightforward and should serve as a warm up for the more involved proof of \Cref{thm: MinDegGen}.
	
	%comment the old formulation of the theorem out.
	
	\iffalse
	\begin{reptheorem}{thm:main}
		Let $k \geq 3$ and $q$ a prime power such that $q \geq 2^{10}k\log q$. For all non-negative integers $s \leq q^2/(2^{12}k\log q)$ there exist edge-disjoint $K_{k+1}$-free graphs $G_1,\dots,G_s$ on the same vertex set $[N]$, $q^4/2 \leq N \leq q^4$, with the additional property that every subset of size at least $2^{100k}q^2\log q$ induces a copy of $K_k$ in each $G_i$.
	\end{reptheorem}
	\fi 
	
	
	\begin{proof}[Proof of \Cref{thm: MinDegConst}]
		Fix some $k \geq 3$, set $C_k=2^{300k}$, and let $r$ be sufficiently large (so that $\sqrt{r}\geq 128\log r$ is satisfied.) By \Cref{obs: ColorPatter}, it suffices to find a $K_{k+1}$\textit{-free} color pattern $G_1,\dots,G_r$ on a vertex set of size $n\leq C_kr^2\log^2r$ such that $\alpha_k(G_i)<n/r$ for all $i \in [r]$. By Chebyshev's theorem, there exists a prime number $q$ satisfying \[\frac{1}{2}\Bigl(C_kr^2\log^2r\Bigr)^{1/4}\leq q \leq \Bigl( C_k r^2 \log^2r \Bigr)^{1/4}.\] Set $c=\lceil \frac{q}{128\log q}\rceil$ and note that 
		\[c\Bigl\lfloor\frac{q}{2}\Bigr \rfloor \geq \frac{q^2}{2^9 \log q} \geq \frac{2^4q^2}{\sqrt{C_k} \log r} \geq r.\]
		Therefore applying \Cref{lem: graphlemma} with the above value of $c$ gives us at least $r$ edge disjoint graphs $\tilde{G}_1,\dots,\tilde{G}_r$ on a shared vertex set of size $\frac{q^4}{2}\leq n\leq q^4=C_k(r \log r)^2$ satisfying $(1)-(5)$ of \Cref{lem: graphlemma} above. By \Cref{Rem:MVConstant}, we can find for each $i \in [r]$ a $K_{k+1}$\textit{-free} subgraph $G_i \subset \tilde{G}_i$ such that
		\[\alpha_{k}(G_i) \leq 2^{40k+8}q^2 \log q \leq 2^{40k+8}\sqrt{C_k} \:r\log r \Bigl(\log r +\log C_k \Bigr)  < \frac{n}{r}. \]
		The graphs $G_1,\dots,G_r$ certify the assertion of \Cref{thm: MinDegConst}.
	\end{proof}
	
	% Comment the proof in the older formulation out.
	\iffalse
	\begin{proof}[Proof of Theorem ~\ref{thm:main}]
		
		
		
		Let $k \geq 3$, $q$ a prime power such that $q \geq 2^{10}k\log q$ and set $c = (q+1)/(2^{10} k \log q)$. Since $c \leq q/(48 \log q)$, we can apply \Cref{lem: graphlemma}  and find $c\lfloor\frac{q}{2}\rfloor$ edge-disjoint graphs $\tilde{G}_i$ on the common vertex set $|L|$. In this setup, it was shown by Mubayi and the second author \cite{mubayi2024order} that we can find a $K_{k+1}$-free subgraph $G_i$ of $\tilde{G}_i$ such that every subset of size at least $2^{100k}q^2\log q$ induces a copy of $K_k$. By construction, all graphs $\tilde{G}_i$ are edge-disjoint and so clearly their subgraphs $G_i$ will be too. We are done as long as $s\leq c\lfloor\frac{q}{2}\rfloor$ which is true since 
		\[s \leq \frac{q^2}{2^{12}k\log q} \leq \frac{q+1}{2^{10}k\log q} \cdot \left\lfloor\frac{q}{2}\right\rfloor . \qedhere\]
	\end{proof}
	
	
	Clearly, by varying the size of $\Lambda$, we can obtain constructions with more colors and fewer vertices or vice versa (see \Cref{remark:morecolors}). {\color{blue} It seems to me that this trade-off is in fact a win-win. I do not understand this remark completely.}
	\textcolor{purple}{SM: I suppose this remark is not really relevant, since the asymptotics will not change. I think the point is that one can take $cq$ of the unitals in pencil, for some constant $c > 0$, and then the number of vertices is $(1-c)q^4$, but we do not care for the constant anyway.}
	Since there is an exponential dependence on $k$ in the previous result, we will now construct a slightly different color pattern and prove \Cref{thm:main2}.\\
	\fi
	
	Note that the use of \Cref{Rem:MVConstant} necessitates the exponential dependency on $k$. To circumvent this by-product we alter our construction slightly; in particular, we modify the random sparsification used in \Cite{mubayi2024order} to prove \Cref{thm: MinDegGen}. The starting point of the proof is again \Cref{lem: graphlemma} and in particular item (4): Every $K_{k+1}$ in a graph $\tilde{G}_i$ is highly structured. For the following proof we set $C:= 2^{100}$.
	
	\begin{proof}[Proof of \Cref{thm: MinDegGen}]
		By \Cref{obs: ColorPatter} it suffices to find a $K_{k+1}$\textit{-free} color pattern $G_1,\dots,G_r$ on a vertex set $L$ of size $n\leq C^4k^2r^{2+\frac{30}{k}}\log^{20}r\log^{20}k$ such that $\alpha_k(G_i) < n/r$ for all $i \in [r]$. To that end, we first use Chebyshev's theorem to find a prime $q$ such that 
		\[\frac{C}{2}k^{\frac{1}{2}}r^{\frac{1}{2}+ \frac{15}{2k}}\log^{5} r \log^{5}k\leq q \leq Ck^{\frac{1}{2}}r^{\frac{1}{2}+ \frac{15}{2k}}\log^{5} r \log^{5}k.\]
		We quickly note the following implication which will be relevant later in the proof: \[\log q \leq \frac{1}{2} \log k + \frac{k+15}{2k} \log r + \mathcal{O}(\log\log r) \leq \log r \log k,\]
		for $r$ and $k$ sufficiently large.
		We choose $c= \lceil \frac{8r}{q}\rceil$ and note that indeed $c \leq \frac{q}{48 \log q}$, and the condition for Lemma \ref{lem:coloringpoints} is satisfied. Since $\lceil \frac{8r}{q}\rceil \lfloor \frac{q}{2}\rfloor \geq r$, applying \Cref{lem: graphlemma} gives us $r$ graphs $\tilde{G}_1,\dots,\tilde{G}_r$ on the vertex set $L$ where $q^4/2\leq |L|\leq q^4$. 
		Our job is done once we find inside every $\tilde{G}_i$ a $K_{k+1}\textit{-free}$ spanning subgraph $G_i$ satisfying $\alpha_k(G_i)<\frac{|L|}{r}$. Without loss of generality we fix the subscript $i$ to $1$ and probabilistically prove that such a subgraph $G_1 \subset \tilde{G}_1$ exists. We fix $\alpha:= r^{-\frac{15}{2k} }\log^{-4} r \log^{-4} k$ and carry out the following random procedure: We vertex partition each point-clique of $\tilde{G}_1$ independently into $k+1$ parts $R_0,\dots,R_k$ where every vertex  is independently placed in $R_j$ with probability $\frac{\alpha}{k}$ for $j \in [k]$ and in $R_0$ with probability $1-\alpha$. We then keep an edge $ab$ in $\tilde{G}_1$ if and only if there is some distinct $i,j \in [k]$ such that $a \in R_i$ and $b \in R_j$ in the partition of the unique point-clique containing $a$ and $b$. Each vertex is placed independently from other vertices and the partitions among the cliques are mutually independent. Let $\hat{G}$ be the probability space corresponding to our random procedure. We make the following two claims whose proofs will be slightly postponed:
		
		\begin{claim}\label{claim: ManyKk}
			\begin{equation}
				\mathbb{P}\Biggl( \exists A \in \binom{L}{|L|/r}: K_k \nsubseteq\hat{G}[A] \Biggr) < \frac{1}{2}.
			\end{equation}
		\end{claim}
		\begin{claim}\label{claim: NoKk+1}
			\begin{equation}
				\mathbb{P} \Biggl( K_{k+1}  \subseteq \hat{G}\Biggr) < \frac{1}{2}. 
			\end{equation}
			
		\end{claim}
		Given the above two claims, we choose for every $i \in [r]$, a $K_{k+1}$\textit{-free} spanning subgraph $G_i \subseteq \tilde{G}_i$ such that $\alpha_k(G_i)<\frac{|L|}{r}$. $G_1,\dots,G_r$ is then the desired color pattern. This proves \Cref{thm: MinDegGen} pending the proofs of the claims which we do next.
		
		
		\begin{proof}[Proof of Claim \ref{claim: ManyKk}]
			Fix some $A \in \binom{L}{|L|/r}$ and let $E_A$ be the event that $K_k \nsubseteq\hat{G}[A]$. Given a point-clique $K$, we say that $K$ is $A$\textit{-proper} if 
			$A_K:= A\cap K$ has size at least $\frac{q^2}{2^4r}=:t$. Furthermore, let us denote $A$\textit{-proper} point-cliques by $P_A$. By a simple double counting we get that:
			\begin{equation}
				\underset{K \in P_A}{\sum}|A_K| \geq \frac{|L|q}{2cr}- \frac{2q^3 t}{c}\geq \frac{|L|q}{4cr}.
			\end{equation}
			%\red{Is the 1/4 in the end correct?}
            %\textcolor{blue}{For this to follow we just need that $\frac{|L|q}{2cr} \geq 2\frac{2q^3t}{c}$ and since $|L|\geq \frac{q^4}{2}$ and $t=\frac{q^2}{2^4r}$ we should be good.} \red{Cool, probably I messed up the powers of 2.}
			For every $A$\textit{-proper} point-clique $K$, let us further vertex partition $A_K$ into $\lfloor \frac{|A_K|}{t}\rfloor=:s_K$ sets $K^1,\dots,K^{s_K}$ each of size $t$. We note that 
			\[E_A\subseteq\{K_k \nsubseteq \hat{G}[K^1],\dots,\hat{G}[K^{s_K}]\} \]
			and that the events $\{K_k\nsubseteq \hat{G}[K^i]\}_{i \in [s_K]}$ are mutually independent. Therefore, we have
			\begin{equation*}
				\mathbb{P}( E_A) = \underset{K \in P_A}{\prod}\underset{j \in [s_K]}{\prod}\mathbb{P}\Bigl(K_k\nsubseteq\hat{G}[K^j]\Bigr) \leq  \underset{K \in P_A}{\prod}\underset{j \in [s_K]}{\prod}k\Bigl(1-\frac{\alpha}{k}\Bigr)^t \leq \exp(-\underset{K \in P_A}\sum\frac{\alpha t}{2k} s_K) \leq \exp\Bigl(-\frac{\alpha }{4k}\underset{K \in P_A}\sum|A_K|\Bigr).
			\end{equation*}
			Where the first inequality follows by the union bound, the second by using $e^{-x}\geq 1-x$ for $x \in (0,1)$ and that $\frac{\alpha t}{k}\geq 2 \log k$, and the last by noting that $\lfloor \frac{|A_K|}{t}\rfloor \geq \frac{|A_K|}{2t}$.
			Using the union bound we can conclude:
			\begin{equation*}
				\mathbb{P}\Bigl( \underset{A \in \binom{L}{|L|/r}}{ \cup}E_A\Bigr)\leq \binom{|L|}{|L|/r} \exp \Bigl( -\frac{\alpha}{4k} \underset{K \in P_A}{\sum}|A_K| \Bigr) \leq \exp \Bigl( \frac{|L| \ln (er)}{r}-\frac{\alpha |L|q}{16kcr}\Bigr) < \frac{1}{2}.
			\end{equation*}
			The last inequality follows since 
            \begin{equation}\label{eq: techKk}
                \alpha q \geq 32kc\ln(er),
            \end{equation}
            
        
			which follows by our choice of $q$ and since $k < r \log^2 r$.
			%\textcolor{blue}{Please decide if this elaboration is needed:\\
            %\underline{Case 1: $c>1$} then $c= \lceil \frac{8r}{q}\rceil \leq \frac{16r}{q} $ and \Cref{eq: techKk} 
            % follows since $q^2\geq 2^9 kr \log^4k \log^4r \ln(er)$ by our choice of $q$. \\
            % \underline{Case 2: $c= 1$}. \Cref{eq: techKk} becomes $q\geq 2^5k \log^4k \log^4r\ln(er)$ which is true by our choice of $q$ and the inequality $k \leq r \log^2 r.$ 
             %}
             %\red{Is the $r^{\frac{15}{2k}}$ from $\alpha$ missing? Anyway, I had no issue with this estimation, so the clarification is not necessary as far as I am concerned.}
		\end{proof}
		\begin{proof}[Proof of Claim \ref{claim: NoKk+1}] Recall that there are only two possible types of $K_{k+1}$ in $\tilde{G}_1$: Degenerate $K_{k+1}$'s which get deleted by the Turánization of the point-cliques; and $K_{k+1}$'s that correspond to $(k+1)$\textit{-fans}. If $P_1$ is the set of point-cliques in $\tilde{G}_1$, then there are at most 
			\[|P_1|.|L| \binom{2q/c}{k} \leq q^7 \binom{2q/c}{k}\]
			$(k+1)$\textit{-fans} in $\tilde{G}_1$.
			Let $F$ be a $(k+1)$\textit{-fan} in $\tilde{G}_1$, and let $K_0,\dots,K_k$ be the point-cliques containing at least two vertices  of $F$, which we shall call relevant. Without loss of generality $|K_0\cap F|=k$  while $|K_i\cap F|=2$ for every $i\in [k]$. The probability that $\hat{G}[F]\cong K_{k+1}$ is at most $\alpha^{3k}$ since every vertex in $F$ must be thrown in the "active portion" of the partition of every relevant point-clique containing it. Therefore, using the union bound again we conclude that:
			\[\mathbb{P}\Bigl(K_{k+1} \subset \hat{G} \Bigr) \leq q^7 \binom{2q/c}{k}\Bigl(\frac{1}{r^{45/2k}\log^{12}k \log^{12}r}\Bigr)^k\leq \exp \Bigl(7\ln q-k\ln\log q -7.4 \ln r\Bigr)<\frac{1}{2}.\]
			for large enough $k$. 
            %\textcolor{blue}{Please decide if the following elaboration is needed.
            The second inequality follows since 
            \[\binom{2q/c}{k} \Bigl(\frac{1}{r^{\frac{45}{2k}}\log^{12}k\log^{12}r}\Bigr)^k \leq \Bigl(\frac{2qe}{ck} \Bigr)^k\Bigl(\frac{1}{r^{\frac{45}{2k}}\log^{12}k\log^{12}r}\Bigr)^k \leq \Bigl(\frac{1}{r^{\frac{15}{2k}}\log r \log k}\Bigr)^k \leq \Bigl(\frac{1}{r^{7.5/k} \log q}\Bigr)^k.\]
            %}
            %\red{This elaboration I do like. However, where is the $C^2$ coming from? Also how does the 7.5 change to a 7.4?
            %}
		\end{proof}

        This completes the proofs of Claims \ref{claim: ManyKk} and \ref{claim: NoKk+1} and hence the proof of \Cref{thm: MinDegGen}.
	\end{proof}
	
	
	
	
	
	\section{Lower Bound}\label{sec:lowerbounds}
	While the best (partial) upper bounds for the minimum degree of minimal Ramsey graphs are quadratic in both parameters (suppressing logarithmic factors), we do not currently have a matching lower bound. We know that $\Omega(rk^2)=\mathrm{ssat}_r(K_k)\leq s_r(K_{k})$ where the equality follows by a theorem of Damásdi et al.\ \cite{damasdi2021saturation} and the inequality by an observation of Tran \cite{tran2022two}. In this section we provide another lower bound quadratic in $r$.
	
	
	%In this section we aim to provide two lower bounds for $s_r(K_{k+1})=P_r(k)$; one quadratic in $r$ and the other quadratic in $k$.\begin{proposition}\label{prop: QuadK}Let $k,r \geq 3$ be positive integers, then $P_r(k) \geq k(r-1)(k-1)+k$.\end{proposition}\begin{proof} Let $G_1,\dots,G_r$ be a $K_{k+1}$\textit{-free} color pattern on the vertex set $V$ such that every vertex coloring of $V$ with colors from $[r]$ contains a strongly monochromatic $K_k$. We shall color $V$ exactly $(r-1)(k-1)+1=:m$ consecutive times $\chi_1,\dots,\chi_m$ and in each time we will find (or "elect") an \textit{$r$-edge-monochromatic} $K_k$ that is vertex disjoint from all the previously \textit{elected} cliques. At each step $i\in [m]$ we maintain vertex-disjoint $r$\textit{-edge-monochromatic} cliques $C_1,\dots,C_{i-1}$, which we call the \textit{elected} cliques, along with a uniform $[r-1]$\textit{-vertex-labeling} of each elected clique, whereas in the first step we have no elected cliques. Uniform here means that any two vertices belonging to the same elected clique must have the same label.  We then color $V$ according to the coloring scheme $\chi_i$ defined as follows: We color the vertices of the cliques $C_1,\dots,C_{i-1}$ according to their label and color everything else with the color $r$.%
		
		
		%As long as the number of cliques $C_1,\dots,C_{i-1}$ whose vertices are labeled with a color $j \neq r$ is at most $k-1$, there will be no strongly monochromatic $K_k$ in color $j$ according to the scheme $\chi_i$. Let $\mathcal{P}$ be the condition that every color $j \in [r-1]$ is used at most $k-1$ times to label cliques. Therefore if $P$ is satisfied, there must exist a strongly monochromatic $K_k$ in color $r$ according to $\chi_i$, which we denote $C_i$. $C_i$ is then vertex disjoint from $C_\ell$ for all $\ell \in [i-1]$. We label all the vertices of $C_i$ by an arbitrary color in $[r-1]$ which was not used at least $k-1$ times to label previous cliques. We can guarantee that condition $\mathcal{P}$ is satisfied for the first $(r-1)(k-1)=m-1$ steps. This together with one last step, in which we give an arbitrary label from $[r-1]$ to our clique, gives us that $|V| \geq \bigl|\underset{i \in [m]}{\dot{\cup}}V(C_i)\bigr|\geq k(k-1)(r-1)+k.$   \end{proof}Notice that the previous argument did not use the fact that the color patterns are $K_{k+1}$\textit{-free}, which potentially signifies the sub-optimality of the bound. In what follows, we shall provide a lower bound on $s_r(K_{k+1})$ that is quadratic in $r$. This has been done in \cite{fox2016minimum}, with a worse factor in $k$.%
	
	\begin{reptheorem}{thm: LowerKk+1}
		For all $k,r \geq 3$, $s_r(K_{k+1}) \geq \frac{kr^2}{16}$.
	\end{reptheorem}
	Let us first prove the following lemma 
	\begin{lemma}\label{lemma: ErdosRogersK}
		Let $k\geq 3$ and let $G$ be a $K_{k+1}$\textit{-free} graph on $n \in \mathbb{N}_{+}$ vertices. There exists a vertex subset $V' \subset V(G)$ such that $|V'| \geq \frac{1}{2}\sqrt{kn}$ and $G[V']$ is $K_{k}$\textit{-free}.
	\end{lemma}
	\begin{proof}
		If a $K_{k+1}$-free graph $G$ has maximum degree $d$, then 
		the neighborhoods of vertices are $K_k$-free. Zykov~\cite{Zykov} proved that the maximum number of copies of $K_{k - 1}$ in a $K_k$-free 
		graph with $d$ vertices is at most $D = (d/(k - 1))^{k - 1}$, with equality achieved by a balanced complete $(k - 1)$-partite graph. So the number of $K_k$ 
		in $G$ is at most $nD/k$. If we randomly sample vertices of the graph with probability $p = D^{-1/(k - 1)}$, then 
		the expected number of vertices remaining after we remove one vertex from each copy of $K_k$ is at least 
		\[ pn - p^k\frac{nD}{k} \geq \Bigl(1 - \frac{1}{k}\Bigr) \frac{n}{D^{1/(k - 1)}} \geq \frac{(k - 1)^2 n}{kd} \geq \frac{kn}{4d}.\] 
		We conclude that $G$ contains a $K_k$-free subgraph $H$ where
		\[ |V(H)| \geq \max\Bigl\{d,\frac{kn}{4d}\Bigr\} \geq \frac{1}{2}\sqrt{kn}.\]
	\end{proof}
	\begin{proof}[Proof of Theorem \ref{thm: LowerKk+1}]
		First, let us note the following recursion for all $r,k\geq 3$.  \begin{equation}\label{recursion}
			P_r(k) \geq P_{r-1}(k)+ \Bigg\lceil \frac{1}{2}\sqrt{kP_{r}(k)}\Bigg\rceil.
		\end{equation}
		
		Indeed, Let $G_1,\dots,G_r$ be an optimal $K_{k+1}$\textit{-free} color pattern on the vertex set $V$ such that every $[r]$\textit{-coloring} of $V$ contains a strongly monochromatic $K_k$. This means $|V|=P_r(k)$ and we can use Lemma \ref{lemma: ErdosRogersK} applied on the $K_{k+1}$\textit{-free} graph $G_r$ to find a subset $V'\subset V$ of size at least $\frac{1}{2}\sqrt{k P_r(k)}$ such that $G_r[V']$ is $K_{k}$\textit{-free}. Let us color $V'$ with the color $r$ and note that every extension of this coloring using colors from $[r-1]$ contains a strongly monochromatic $K_k$ in $V-V'$ in one of the colors $[r-1]$ and therefore, $|V-V'| \geq P_{r-1}(k)$ proving our claim. Next, we claim that for $k,r \geq 3$, we have $P_r(k) \geq \frac{kr^2}{16}$. Note that this claim is equivalent to the statement of \Cref{thm: LowerKk+1} by \Cref{lem: ColorPattern}. We proceed by induction on $r$. Fix some $k \geq 3$, for the base case of our induction we have $r=3$. However, $P_3(k) \geq P_2(k)=k^2\geq \frac{3^2k}{16}.$ Next, we employ the recursion in \ref{recursion} to have \begin{equation}
			P_r(k)-\frac{1}{2}\sqrt{kP_r(k)} \geq P_{r-1}(k) \geq \frac{k(r-1)^2}{16}, 
		\end{equation}
		where the last inequality follows by the induction hypothesis. Let $x:= \sqrt{P_r(k)}$, then we know that
		\[x^2-\frac{1}{2}\sqrt{k}x - \frac{k(r-1)^2}{16} \geq 0\]
		or equivalently 
		\[\Biggl(x-\frac{\sqrt{k}}{4} -\sqrt{\frac{k}{16}+\frac{k(r-1)^2}{16}}\Biggr) \Biggl(x-\frac{\sqrt{k}}{4} +\sqrt{\frac{k}{16}+\frac{k(r-1)^2}{16}} \Biggr)\geq 0.\]
		Then either $\sqrt{P_r(k)} =x\leq \frac{\sqrt{k}}{4}-\sqrt{\frac{k}{16}+\frac{k(r-1)^2}{16}}<0$
		which gives us a contradiction; or
		\[\sqrt{P_r(k)}=x \geq \frac{\sqrt{k}}{4} + \sqrt{\frac{k}{16}+\frac{k(r-1)^2}{16}} \geq \sqrt{\frac{kr^2}{16}}\]
		and we are done. 
		%\red{We might have to add another power of 2 in the denominator since $(r-1)^2 \geq r^2$ will not be true ($r$ could be huge and $k$ very small in the above).}
        %\textcolor{blue}{I still think it is correct. the difference between $r$ and $\sqrt{(r-1)^2}$ is approaching a constant as $r$ goes to infinity. This constant (multiplied by a k ) will be accounted for by the other terms in $k$. Here is a formal argument for it: Equivalently by squaring and dividing by $16$, we want to show that 
        %\[2k+ k(r-1)^2+2k\sqrt{1+(r-1)^2}\geq kr^2\] or
        %\[2+ (r-1)^2+2\sqrt{1+(r-1)^2}\geq r^2\] or equivalently 
        %\[r\leq 3/2+\sqrt{1+(r-1)^2}\] and this is obviously true.}
        %\red{Ok cool, I did not take the whole thing into account. I wonder if the average reader will, but we can leave it as is.}
	\end{proof}
	
	
	
	
	\section{Semisaturated Ramsey Numbers}\label{section: Tran}
	In this short section we prove Theorem~\ref{thm:semisaturation} about our upper bound on semisaturated Ramsey numbers. 
    %As mentioned in the introduction, Tran already observed that $\mathrm{ssat}_r(K_{k+1}) \leq P_r(k)$. 
   % The next theorem provides an upper bound on $\mathrm{ssat}_r(K_k)$.
%	
%	\begin{reptheorem}{thm:semisaturation}
%		For all $k,r \geq 2$, we have $\mathrm{ssat}_r(K_{k+1}) \leq 4k^2r^2$.
%	\end{reptheorem}
	\iffalse
	\begin{proposition}\label{prop:saturationinequalities}
		For all $k,r \geq 2$, we have $\mathrm{ssat}_r(K_{k+1}) \leq P_r(k) \leq \mathrm{sat}_r(K_{k+1})$.	
	\end{proposition}	
	
	\begin{proof}
		For the first inequality, one can take a $K_{k+1}$-free color pattern $G_1,\dots,G_r$ on $N = P_r(k)$ vertices and color the remaining edges of $K_N$ (i.e.\ those not covered by any $G_i$) arbitrarily to obtain an edge-colored graph $G \in \cG_r$. Now consider $G' \in \cG_r$ that extends $G$ and has one vertex $v$ more. Define a vertex-coloring $\phi$ of $G$ by setting $\phi(w)$, $w \in V(G)$, to be the color of the edge $vw$ in $G'$. By definition of $P_r(k)$, we then obtain a strongly monochromatic copy of $K_k$, which in $G'$ corresponds to a copy of $K_{k+1}$ by adding the vertex $v$. 
		
		For the second inequality, consider the smallest $(r,K_{k+1})$-saturated graph $G$ on $N = \mathrm{sat}_r(K_{k+1})$ vertices. This coloring defines a $K_{k+1}$-free color pattern $G_1,\dots,G_r$ on a vertex set $V$ of size $N$ by setting $G_i$ to be the graph induced by color $i$. Moreover, any coloring $\phi$ of $V$ will induce a strongly monochromatic $K_k$: we can translate this coloring to an edge-coloring of $K_{N+1}$ by adding one extra vertex $v$ and defining the color of the edge $vw$ to be $\phi(w)$ and all other edges receive the same color as in $G$. By definition of $\mathrm{sat}_r(K_{k+1})$, this coloring will produce a copy of $K_{k+1}$ in some color $i$, which implies that color $i$ induces a strongly monochromatic copy of $K_k$ in the color pattern.
	\end{proof}
	\fi
	\begin{proof} For our proof we must construct an edge $r$-coloring of $K_n$ with $n=4(k-1)^2r^2$ such that any extension of it to an $r$-edge coloring of $K_{n+1}$ creates a new monochromatic $K_{k+1}$. 

		By Chebyshev's theorem, we can find a prime $q$ such that $(k-1)r < q < 2(k-1)r$.         
        Consider the affine plane $(\cP,\cL)$ of order $q$ and denote its $q+1$ parallel classes of lines by $\cL_1,\dots,\cL_{q+1}$. 
        Define first an edge $(q+1)$-coloring of the complete graph $K_{q^2}$
        with vertex set $\cP$ as follows: 
   %     Define a graph $G \in \cG_{q+1}$ on the points of the affine plane as follows: 
   the edge between two vertices receives color $i$ if and only if the line they span on the affine plane is in $\cL_i$. One can observe that every color class is then the union of $q$ disjoint cliques of order $q$, and any two cliques of a different color meet in exactly one vertex. To create an $r$-coloring we can for example collapse the color classes between $r$ and $q+1$ into one.\\
	We need to show that adding a vertex and $r$-coloring the edges incident to it necessarily creates a new monochromatic $K_{k+1}$. There is certainly a color, say color $i \in [r]$, that occurs at least the average number, i.e. $q^2/r > (k-1)q$ times among the $q^2$ edges incident to the new vertex. Since there are $q$ pairwise disjoint monochromatic cliques in color $i$ covering the whole vertex set, one of these cliques must have at least $k$ vertices which are connected to the new vertex via an edge of color $i$. These vertices, together with the new vertex, form a new $K_{k+1}$ that is monochromatic in color $i$.    
 %   By passing to the largest color class in the neighborhood of the added vertex, it suffices to show that for all $i \in [q+1]$ every set of size at least $|\cP|/r$ induces a copy of $K_k$ in color $i$. This will prove the theorem,
%		as the number of vertices of $G$ is $q^2 < 4(k-1)^2r^2 < 4k^2r^2$ and one can easily recolor the graph $G$ to obtain a member $G' \in \cG_r$ by taking unions of color classes. Formally: for any surjective function $f:[q+1]\to[r]$, we can define a graph $G' \in \cG_r$ by defining $G_j' = \cup_i G_i$, where the index $i$ runs over all $i \in [q+1]$ such that $f(i) = j$, and $G_j'$ (resp.\ $G_i$) denotes the subgraph of $G'$ (resp.\ of $G$) induced by color $j$ (resp. $i$).   \\
%		
%		So fix a color $i \in [q+1]$ and consider a set of points $U$ of size at least $|\cP|/r = q^2/r > (k-1)q$. Since there are exactly $q$ lines in a parallel class, we find by pigeonhole principle that there must be a line $\ell \in \cL_i$ containing $k$ points of $U$. By definition of $G$, this corresponds to a copy of $K_k$ in color $i$.
	\end{proof}
	
	\section{Concluding remarks} \label{sec:remarks}
	
	$\bullet$ {\bf Asymptotics of $s_r(K_{k})$.} 
	Tantalizing problems remain open in all ranges of the parameters. 
	When $k$ and $r=r(k)$ both tend to infinity, the main question concerns correct exponents of $k$ and $r$. 
	If $k\leq r \log^2 r$, Theorems~\ref{thm: MinDegGen} and \ref{thm: LowerKk+1} give 
	$$ \frac{1}{16}r^2 k \leq s_r(K_{k})  \leq (rk)^{2+o(1)},$$
	so the exponent of $k$ is waiting to be settled.
	If $k > r \log^2 r$, then the bound of H\`an et al.\ ~\cite{han} gives 
	$$r k^2(1+o(1)) \leq s_r(K_{k})  \leq r^{3+o(1)} k^{2+o(1)},$$
	which leaves the exponent of $r$ up for grabs.
	
	For constant $k\geq 4$, Theorem~\ref{thm: MinDegConst} and the lower bound (\ref{eq:fox-constantk}) of Fox et al.\ ~\cite{fox2016minimum} gives
	$$ c_k r^2 \frac{\log r}{\log\log r} \leq s_r(K_{k})  \leq C_k r^2 \log^2 r,$$
	so the status of a factor of $\log r \log \log r$ remains unsettled. 
	For constant $r\geq 3$, the bounds 
	$$ c_r k^2 \leq s_r(K_{k})  \leq C_rk^2 \log^2 k.$$ of H\`an et al.\ \cite{han} are leaving us with the challenge whether how much of the factor $\log^2 k$ is necessary. 
	Making progress on any of these bounds would be very interesting. 
	We are especially eager to discover new avenues to obtain lower bounds for the problem. 
	
	For completeness we recall that for $r=2$ colors the exact value $s_2(k)=(k-1)^2$ was determined by Burr et al.\ ~\cite{burr1976graphs} and for $k=3$ the order of magnitude $s_r(3) = \Theta (r^2\log r)$ is known by Guo and Warnke~\cite{guo2020packing} and Fox et al.\ \cite{fox2016minimum}. 
    The determination of the constant factor in the latter problem is related to the analogous question for the Ramsey number $R(3,\ell)$, a notorious open problem. 
	
	
	$\bullet$ {\bf Separation in the (semi)saturation problem.} 

By the definition of the $r$-color Ramsey number $R_r(K_{k+1})$, there exists an edge $r$-coloring of the complete graph on $R_r(K_{k+1})-1$ vertices, which does not contain any monochromatic $K_{k+1}$. Moreover, $R_r(K_{k+1}) -1$ is the largest number of vertices on which such an edge $r$-coloring exists and hence this coloring is in the family ${\cal RC}_r(K_{k+1})$.  
%	Let ${\cal RC}_r(K_k)$ be the set of $r$-colorings of some complete graph, such that any extension of it to an $r$-coloring of the edges of the complete graph with one more vertex creates a new monochromatic copy of $K_k$. The motivation behind this notion is that the classical $r$-color Ramsey number $R_r(K_k)$ is exactly $1$ plus the {\em largest} integer $n$ such that there exists an $r$-coloring of $E(K_n)$ in ${\cal RC}_r(K_k)$ without a monochromatic $K_k$. 
	Dam\'asdi et al.\ ~\cite{damasdi2021saturation} defined $\mathrm{sat}_r(K_{k+1})$ to be the {\em smallest} integer $n$ such that there exists an $r$-coloring of $E(K_n)$ in ${\cal RC}_r(K_{k+1})$ without any monochromatic $K_{k+1}$. Thus $\mathrm{sat}_r(K_{k+1}) < R_r(K_{k+1})$. 
	This concept was inspired by the definition of Erd\H os, Hajnal and Moon~\cite{erdos1964problem} of the saturation number $\mathrm{sat}(n,H)$ of graph $H$ opposite of its Tur\'an number $\mathrm{ex}(n,H)$. 
 
 Then trivially $\mathrm{ssat}_r(K_{k+1}) \leq \mathrm{sat}_r(K_{k+1})$, as in the definition of $\mathrm{ssat}_r(K_{k+1})$ we are allowed to consider {\em all} members of ${\cal RC}_r(K_{k+1})$, not only those without monochromatic $K_{k+1}$. We can also easily see that our main focus, the smallest minimum degree $s_r(K_{k+1}) = P_r(k)$ is sandwiched between these two saturation parameters. 
	
	\begin{proposition} \label{prop:ssat-P-sat}
		$\mathrm{ssat}_r(K_{k+1}) \leq P_r(k) \leq \mathrm{sat}_r(K_{k+1})$
	\end{proposition}
	\begin{proof}
		 For the second inequality, we connect the nomenclature of $P_r(k)$ to the color saturation parameter $\mathrm{sat}_r(K_{k+1})$. Firstly, an edge $r$-coloring of $K_n$ can be identified with an $r$-color pattern $G_1, \ldots , G_r$ on $[n]$, such that $\cup_{i=1}^r E(G_i) = E(K_n)$. Moreover, an edge $r$-coloring having the property that every extension of it to $n+1$ vertices creates a new monochromatic $K_{k+1}$ is equivalent to the corresponding $r$-color pattern $G_1, \ldots , G_r$ having the property that every $r$-coloring of $[n]$ contains a strongly monochromatic $K_k$. Indeed, extensions of an edge $r$-coloring of $K_n$ to the edges incident to the new vertex $(n+1)$ are in one-to-one correspondence with the vertex $r$-colorings of $[n]$ (namely, the color of an extension edge is just the color of the endpoint of that edge in $[n]$ in the vertex coloring). Then having a ''new'' monochromatic $K_{k+1}$ in the extension is equivalent to having a strongly monochromatic $K_k$ in the vertex coloring. 
		The second inequality of the proposition then follows because for the definition of $P_r(k)$ we do {\em not} require the $r$-color pattern to partition the whole edge set of the clique, while for $\mathrm{sat}_r(K_{k+1})$ we do, so the minimum $n$ for the latter is taken over a subset of the set we take the minimum of for the former. 
        
        For the first inequality, which has already been observed by Tran~\cite{tran2022two}, one can take an $r$-color pattern $G_1, \ldots , G_r$ on the optimal number $n=P_r(k)$ of vertices and color arbitrarily the uncolored edges in $E(K_n) \setminus \cup E(G_i)$ by $r$ colors. This provides an $r$-coloring of $E(K_n)$ that is in ${\cal RC}_r(K_{k+1})$. 
     \end{proof}

It is natural to wonder how tight the two inequalities of the previous proposition are in the various ranges of the parameters. 

For $r=2$ for example it is known \cite{damasdi2021saturation} that all three functions are equal to $k^2$. 
    %both inequalities of the previous proposition are equalities.
On the other hand for any constant $k\geq 2$ our Theorem~\ref{thm:semisaturation} and the lower bounds of \cite{fox2016minimum} do separate the order of magnitude of $\mathrm{ssat}_r(K_{k+1})$ and $P_r(k)$ as $r$ tends to infinity. 

We believe that this to be true for any other ranges of the parameters.

\begin{conjecture} \label{con:ssat-P} For any $r=r(k) \geq 3$, as $k$ tends to infinity we have 
		$\mathrm{ssat}_r(K_{k+1})\ll P_r(k)$. 
        \end{conjecture}
    
   	More modestly, it would even be interesting to decide whether there is significant separation between $\mathrm{ssat}_r(K_{k+1})$ and $\mathrm{sat}_r(K_{k+1})$.  
 
For $r=3$ H\`an, R\"odl and Szab\'o~\cite[Conjecture 2]{han}  conjectured an asymptotic separation between the three- and two-color case of the smallest minimum degree parameter. This is plausible, yet we do not even know whether there is {\em any} $r$ for which 
$s_r(K_{k+1}) = P_r(k) \gg s_2(K_{k+1}) = k^2$ as $k \to \infty$. 
Since $P_r(k)$ is non-decreasing in $r$ and by Theorem~\ref{thm:semisaturation} the order of $\mathrm{ssat}_r(K_{k+1})$ is quadratic in $k$ for every fixed $r$, the conjecture of H\`an et al.\ would also imply Conjecture~\ref{con:ssat-P} for any fixed $r\geq 3$.

    Concerning the separation in the second inequality of Proposition~\ref{prop:ssat-P-sat} we are less convinced. For the case of $k=2$ we tend to agree with authors of \cite{guo2020packing} who believe that their construction could be improved so that all (and not only a $(1-\varepsilon)$ proportion) of the edges of the complete graph are covered by some $K_3$-free $r$-color pattern on $\Theta ( r^2\log r)$ vertices. 

\begin{conjecture}
    $P_r(2) = \Theta (\mathrm{sat}_r(K_{3})).$ 
    \end{conjecture}


    
	$\bullet$ {\bf Monotonicity of $s_r(K_k)$ in $k$.} Finally, we would like to reiterate the humbling conjecture of Fox et al.~\cite[Conjecture 5.1]{fox2016minimum} stating that $s_r(K_{k+1}) \geq s_r(K_k)$ for every $r\geq 2$ and $k\geq 2$. Recall that from the identity $s_r(K_{k+1})=P_r(k)$, it is not difficult to see that $s_r(K_{k+1})$ is non-decreasing in $r$. 
	
	
	
	%$\bullet$ {\bf Lower bounds on the Erd\H{o}s-Rogers function.} Given two positive integers $k$ and $n$, the Erd\H{o}s-Rogers function $f_k(n)$ is the largest integer $m$ such that every $K_{k+1}$\textit{-free} graph on $n$ vertices contains a $K_k$\textit{-free} vertex subset of size $m$. Lower bounds on $f_k(n)$ lead to lower bounds on $P_r(k)$ through the following recursion :
	%\[P_{r}(k) \geq P_{r-1}(k)+ f_k(P_r(k)).\]
	%Indeed, this approach was used in \Cref{sec:lowerbounds}. It is of interest then to provide lower bounds on the Erd\H{o}s-Rogers function $f_k(n)$. The best known lower bound is due to Dudek and Mubayi and reads as follows: 
	%\[f_k(n)= \Omega\Bigl(\frac{\sqrt{n\log n} }{\log \log n}\Bigr).\]
	%The proof proceeds as follows: A $K_{k+1}$ graph $G$ with maximum degree $\Delta$ contains a $K_k$\textit{-free} subset of size $\Delta$ and by a result of Shearer ~\cite{shearer1995independence} an independent set of size $\Omega \bigl( \frac{n\log \Delta}{ \Delta \log \log \Delta }\bigr)$. This $1/\log \log$ factor in the lower bound of $f_k(n)$ is inherited from Shearer's result. Shearer's result however gives a large independent set, a much stronger condition than $K_k$\textit{-freeness}. Is it possible to provide better lower bounds on the Erd\H{o}s-Rogers function by appropriating Shearer's Lemma or abandoning it all together? \textcolor{blue}{I just realized that this has been asked/mentioned in the Mubayi-Verstraete paper.}\\
	
	
	
	
	
	
	
	
	\printbibliography
	
	
\end{document}




